\section{Shared results I: Hodge--Tate, Fontaine and comparison data}\label{sec:pf:sh-hodge}

\providecommand{\shhodgeinv}{\operatorname{inv}}
\providecommand{\shhodgeHf}{H^1_{f}}
\providecommand{\shhodgeZf}{Z^1_{f}}
\providecommand{\shhodgeHfD}{H^1_{f,\mathrm{div}}}
\providecommand{\shhodgecl}{\operatorname{cl}}
\providecommand{\shhodgelen}{\operatorname{len}}
\providecommand{\shhodgeDec}{D}
\providecommand{\shhodgeIne}{I}
\providecommand{\shhodgeNm}{\mathrm{N}}
\providecommand{\shhodgeTr}{\operatorname{Tr}}
\providecommand{\shhodgeres}{\operatorname{res}}
\providecommand{\shhodgeSel}{\operatorname{Sel}}
\providecommand{\shhodgeBK}{\textup{(BK)}}
\providecommand{\shhodgeLTD}{\textup{(LTD)}}
\providecommand{\shhodgeLEC}{\textup{(LEC)}}
\providecommand{\shhodgeGAR}{\textup{(GAR)}}

\subsection{Introduction}\label{pf:sh-hodge:intro}

This section proves sixteen results of the formal development \cite{LeanSource2026} which are used by several
sections of Part~II and which concern the Galois-theoretic side of the proof: local Galois cohomology, the
Hodge--Tate, Fontaine and comparison data of a datum, Wach modules, and three global statements. They are grouped as
follows.

\begin{enumerate}[label=(\roman*)]
\item \emph{Two classical statements} (\cref{pf:sh-hodge:classical}): a degree-one form of Chebotarev's density
  theorem (\cref{pf:sh-hodge:cheb}), used in \cref{sec:pf:galred,sec:pf:kato}; and the deduction of Poitou--Tate
  duality, in the form used by the development, from its \emph{normalised} form (\cref{pf:sh-hodge:pt}), used in
  \cref{sec:pf:liftinj,sec:pf:modularity}.
\item \emph{Local Galois cohomology at all places} (\cref{pf:sh-hodge:bk}): the five Bloch--Kato statements with
  arbitrary coefficients at every prime (\cref{pf:sh-hodge:bkall}), used in \cref{sec:pf:liftinj}, two corollaries
  used in \cref{sec:pf:linkN,sec:pf:linkN-adm,sec:pf:linkN-sel}, and the finiteness of local cohomology of lattices
  (\cref{pf:sh-hodge:fin}).
\item \emph{Hodge--Tate, Fontaine and comparison data} (\cref{pf:sh-hodge:hw}): their existence from the theorem of
  Bhatt--Scholze on crystalline lattices (\cref{pf:sh-hodge:exhw}), used in \cref{sec:pf:linkN}; Herr's description of
  lattice cohomology (\cref{pf:sh-hodge:herr}); and the \emph{later clauses of the Wach lattice theory}
  (\cref{pf:sh-hodge:wach,pf:sh-hodge:used}), used in \cref{sec:pf:linkN-adm,sec:pf:linkN-sel}.
\item \emph{Local torsion bounds} (\cref{pf:sh-hodge:loc}): a bound for the torsion of invariants
  (\cref{pf:sh-hodge:tek}), the twisted torsion class theorem~TB (\cref{pf:sh-hodge:tb}) and the local lifting theory with
  constants $(10,23,32,32)$ (\cref{pf:sh-hodge:lc}), used in \cref{sec:pf:linkN,sec:pf:linkN-adm,sec:pf:linkN-sel}.
\item \emph{Three global statements} (\cref{pf:sh-hodge:glob}): the polynomial bound for the parameters of a family
  (\cref{pf:sh-hodge:par}), the signature choice condition (\cref{pf:sh-hodge:js}) and the global lattice index theorem
  (\cref{pf:sh-hodge:glit}), used in \cref{sec:pf:linkN,sec:pf:linkN-sel}.
\end{enumerate}

\paragraph{The literature hypotheses.}
Most results are stated under some of the following hypotheses, which are propositions of the development
recording published theorems; they are among the components of $\sfname{Hypotheses}_2$ and of the package $P$ of
\cref{sec:families:galois}.
\begin{itemize}
\item[\shhodgeLTD] \emph{Local Tate duality} \cite{Tate1963}, \cite[Ch.~II, \S5]{Serre1997}: for every prime $\ell$, every
  finite extension $K$ of $\Ql$ inside a fixed algebraic closure, every $r\ge1$ and every finite discrete continuous
  $\Z/\ell^r[G_K]$-module $M$, the group $H^2(G_K,\Z/\ell^r(1))$ is isomorphic to $\Z/\ell^r$ and the cup products
  $H^i(G_K,M)\times H^{2-i}(G_K,M^\vee(1))\to H^2(G_K,\Z/\ell^r(1))$, $i=0,1,2$, are perfect.
  \formal{Formalism.GaloisCohomology.LocalTateDualityK}
\item[\shhodgeLEC] \emph{The local Euler characteristic formula} \cite{Tate1963}: for the same $K$ and $M$,
  $\#H^1(G_K,M)=\#H^0(G_K,M)\cdot\#H^2(G_K,M)\cdot(\#M)^{[K:\Ql]}$.
  \formal{Formalism.GaloisCohomology.LocalEulerCharK}
\item[\shhodgeBK] \emph{The local theory of Bloch--Kato} \cite{BlochKato1990}: for every prime $\ell$, every local theory at
  $\ell$ with $A_{\inf}$-data which is \emph{tied} to Fontaine's period rings (\cref{pf:sh-hodge:def-loc}) and every finite
  extension $K/\Ql$: the dimension formula for $H^1_f$ of de~Rham $\Ql$-representations, the rational duality of $H^1_f$,
  and its integral form for $\Zl$-lattices (\cref{pf:sh-hodge:def-bk}).
  \formal{GlobalTheory.Published.BlochKatoLocal}
\item[\shhodgeGAR] \emph{Global Artin reciprocity} (component of $\sfname{Hypotheses}_2$).
  \formal{GlobalTheory.Published.GlobalArtinReciprocity}
\end{itemize}
Two further published statements appear: the theorem of Bhatt--Scholze on crystalline lattices and prismatic
$F$-crystals \cite{BhattScholze2023} (\cref{pf:sh-hodge:def-bs}), and a Chebotarev statement for open normal subgroups
(\cref{pf:sh-hodge:js}).

\paragraph{What is proved here and what is cited.}
The arguments of \cref{pf:sh-hodge:classical,pf:sh-hodge:bk} and of the parameter bound are given in full, except for an
analytic lemma (a Dirichlet theorem for primes of degree one, \cref{pf:sh-hodge:dirichlet}) and the auxiliary lemmas of the
descent from $\Zl$ to $\OO$ and of the finite-level duality, which are stated where they are used. The other results are large
assemblies: the formal proof of the global lattice index theorem alone has about $34\,000$ lines, the later clauses of the Wach
lattice theory about $12\,000$, theorem~TB about $9\,300$, the existence of the Hodge--Tate, Fontaine and comparison data about
$9\,000$ and the local lifting theory about $8\,000$. For these we state the results precisely, give the reduction to their
constituents in full, prove the constituents with short proofs (among them four of the clauses of the Wach lattice theory and the
bound for the torsion of invariants), and state the others, whose proofs are deferred. Peculiarities of the formal statements are
recorded in remarks and collected in \cref{pf:sh-hodge:remarks}.

% ---------------------------------------------------------------------------
\subsection{Conventions}\label{pf:sh-hodge:conv}

We use the notation of \cref{sec:pf:linkN-adm,sec:pf:linkN-sel,sec:pf:linkN} for data. A datum $D$ has rank parameters,
among them $f,n,k,m,r,t$ and $n_V=r+t$, a host with a CM field $E$ of degree $2f$, class number $h_E$, primes $q,d$, type
primes $g_0,g_1$, and a set of local indices $\iota$, each with a prime $\ell_\iota$, a complete discrete valuation ring
$\OO_\iota$ finite over $\Z_{\ell_\iota}$ with fraction field $K_\iota$, eigensystems $X_\iota\supseteq Z_\iota\supseteq
Z^{\mathrm{good}}_\iota$ and finite sets $S_\iota,Q_\iota$ of auxiliary places (\cref{sec:pf:linkN}). Normalised lengths are
$\lambda_\iota(N)=\len_{\OO_\iota}(N)/\len_{\OO_\iota}(\OO_\iota/\ell_\iota)$ ($0$ if $N$ has infinite length).

\begin{definition}[Local theories]\label{pf:sh-hodge:def-loc}
Fix a prime $\ell$, an algebraic closure $\Omega_\ell$ of $\Ql$ and write $G_{\Ql}$ for its automorphism group; finite
extensions $K,\mathit{LL}$ of $\Ql$ are subfields of $\Omega_\ell$, $K_0$ is the maximal unramified subextension of $K$,
$G_K\subseteq G_{\Ql}$ the subgroup fixing $K$, $\OO_{\mathit{LL}}$ the integral closure of $\Zl$ in $\mathit{LL}$, and
$K_s=K_0(\mu_{\ell^s})$.
\begin{enumerate}[label=(\alph*)]
\item A \emph{local theory} at $\ell$ consists of rings $B$, $B_{\mathrm{dR}}$ with actions of $G_{\Ql}$, a Frobenius
  $\varphi$ of $B$, an injection $B\to B_{\mathrm{dR}}$, an element $t\in B$, a character $\chi\colon G_{\Ql}\to\Ql^\times$,
  a compatible system $(\zeta_{\ell^j})_j$ of $\ell$-power roots of unity in $\Omega_\ell$, valuations $v_{\mathit{LL}}$ on
  every $\mathit{LL}$ with $v_{\mathit{LL}}(\ell)=1$ and $\OO_{\mathit{LL}}=\{v_{\mathit{LL}}\ge0\}$, and for every $K$ a period datum
  $\mathrm{PD}_K$ for $G_K$ with these rings (as in \cref{sec:pf:linkN-adm}, \S\ref{pf:linkN-adm:local}); $\chi$ is
  cyclotomic: $g(\zeta)=\zeta^a$ for $\zeta^{\ell^j}=1$ and an integer $a$ with $|\chi(g)-a|\le\ell^{-j}$.
  \formal{Formalism.BasicNotions.LocPrime}
\item $A_{\inf}$-\emph{data} consist of an action of $G_{\Ql}$ on $\C_\ell$ and on the tilt $\OO_{\C_\ell^\flat}$, the element
  $\eps=(\zeta_{\ell^j})_j$ of the tilt, and an injective $G_{\Ql}$- and Frobenius-equivariant map
  $\iota_A\colon A_{\inf}=W(\OO_{\C_\ell^\flat})\to B$ with the usual compatibilities with $\mathrm{Fil}^0$, $\mathrm{Fil}^1$ and
  Teichm\"uller representatives of roots of unity of order prime to $\ell$.
  \formal{Formalism.BasicNotions.LocAinf}
\item The theory with its $A_{\inf}$-data is \emph{tied} if there are ring isomorphisms $B\cong B_{\mathrm{cris}}$,
  $B_{\mathrm{dR}}\cong B_{\mathrm{dR}}$ onto Fontaine's rings \cite{Fontaine1982,Fontaine1994} compatible with the inclusions
  of $\Zl$, the Galois actions, Frobenius, the filtrations, $\theta$, the element $t=\log[\eps]$, and $\iota_A$.
  \formal{Formalism.BasicNotions.LocTied}
\end{enumerate}
Every datum $D$ has a local theory $\mathrm{PD}^D$ at every prime, with $A_{\inf}$-data, and is tied
(\cite{LeanSource2026}, fields of the datum). For a place $w$ of $E$ above $\ell$ the completion $E_w$ is identified with a
finite extension $K_w$ of $\Ql$ and the decomposition group with $G_{K_w}$; the period datum at $w$ is $\mathrm{PD}^D_{K_w}$.
\end{definition}

\begin{definition}[Coefficients, lattices and cohomology]\label{pf:sh-hodge:def-coh}
Let $\mathrm{PD}$ be a period datum for $G=G_K$. A \emph{good coefficient ring} is a pair $(\mathit{LL},\OO)$ with $\mathit{LL}$ a
finite extension of $\Ql$ (with its module topology) and $\OO\subseteq\mathit{LL}$ an open discrete valuation ring with fraction
field $\mathit{LL}$; for every finite $\mathit{LL}/\Ql$ the pair $(\mathit{LL},\OO_{\mathit{LL}})$ is good
(\lref{GlobalTheory.T.padicIntegers_finiteExt_goodCoefficientRing}). A \emph{representation} is a finite-dimensional
$\mathit{LL}$-vector space $V$ with a continuous $\mathit{LL}$-linear action of $G$, and a \emph{lattice} of $V$ is a $G$-stable
finitely generated free $\OO$-submodule $T$ spanning $V$, as in \cref{pf:linkN-adm:def-col}. We write $Z^1(T)$, $B^1(T)$ and
$H^1(T)=Z^1(T)/B^1(T)$ for the continuous cocycles with values in $T$, the coboundaries and the \emph{lattice cohomology};
$\shhodgeZf(V)$ for the \emph{finite} cocycles $c\colon G\to V$ (those with $c(g)\otimes1=gy-y$ for some $y\in V\otimes B$)
\cite{BlochKato1990}, $\shhodgeHf(V)$ for their classes, and $\shhodgeZf(T)=Z^1(T)\cap\shhodgeZf(V)$. For the discrete module
$V/T$ we write $Z^1(V/T)$ for the locally constant cocycles, $B^1(V/T)$ for the coboundaries, and
\[
\shhodgeHfD(T)=\{\bar c:\ c\in\shhodgeZf(V)\}+B^1(V/T)\subseteq Z^1(V/T)
\]
for the module of \emph{finite cocycles of $V/T$} (images of finite cocycles of $V$). For $\varpi\in\OO$ a uniformiser and
$r\ge1$, the finite cocycles of $T/\varpi^rT$ are the reductions of the elements of $\shhodgeZf(T)$, plus coboundaries.
\formal{Formalism.WeaklyAdmissibleFiltrations.HfD}
\formal{Formalism.WeaklyAdmissibleFiltrations.HfQ}
\end{definition}

% ---------------------------------------------------------------------------
\subsection{Two classical statements}\label{pf:sh-hodge:classical}

\subsubsection{Chebotarev's theorem for places of degree one}

For a number field $F$ with algebraic closure $\Omega$ and a finite place $v$ of $F$, a \emph{place of $\Omega$ above $v$} is a
valuation ring $B$ of $\Omega$ whose maximal ideal meets $\OO_F$ in $v$. We write $\shhodgeDec_B\subseteq G_F$ for its
decomposition group, $\shhodgeIne_B$ for the inertia group (the $\sigma\in\shhodgeDec_B$ with $\sigma x\equiv x$ modulo the
maximal ideal for all $x\in B$), and say that $\sigma\in G_F$ is a \emph{Frobenius at $B$ for $q$} if $\sigma\in\shhodgeDec_B$
and $\sigma x\equiv x^q$ modulo the maximal ideal of $B$ for all $x\in B$; $\shhodgeNm v=\#\OO_F/v$.

\begin{theorem}[Chebotarev, degree one]\label{pf:sh-hodge:cheb}
Let $F$ be a number field with algebraic closure $\Omega$, $L\subseteq\Omega$ a finite Galois extension of $F$,
$g\in\Gal(L/F)$ and $S$ a finite set of finite places of $F$. There are a finite place $v\notin S$ of $F$ of absolute residue
degree $1$, a place $B$ of $\Omega$ above $v$ with $\shhodgeIne_B$ fixing $L$, and $\sigma\in G_F$ which is a Frobenius at $B$ for
$\shhodgeNm v$ and restricts to $g$ on $L$.
\formallabel{GlobalTheory.Published.chebotarevDegreeOne'}
\end{theorem}

\begin{remark}[Relation to the literature]\label{pf:sh-hodge:chebrem}
This is a consequence of Chebotarev's density theorem \cite{Chebotarev1926}, \cite[Ch.~VII, Thm.~13.4]{Neukirch1999}: the
places of $F$ unramified in $L$ whose Frobenius class is the class of $g$ have Dirichlet density $\#\mathrm{cl}(g)/[L:F]>0$, and
the places of residue degree $\ge2$ have density $0$, so infinitely many places of degree one have Frobenius class
$\mathrm{cl}(g)$; choosing $B$ and conjugating the Frobenius gives the statement. The formal statement is weaker than the
density theorem (it asserts existence outside a finite set), and the formal proof does not use the density theorem: it
follows Chebotarev's original reduction to cyclotomic extensions (``crossing'' with a cyclotomic field) and proves the
cyclotomic case by the analytic method of Dirichlet and Hecke. We give this proof.
\end{remark}

\begin{lemma}\label{pf:sh-hodge:primem}
For a number field $K$ and $n\ge1$ there is a prime $m$ with $n\mid m-1$ such that the cyclotomic polynomial $\Phi_m$ is
irreducible over $K$.
\formallabel{GlobalTheory.T.exists_prime_dvd_sub_one_cyclotomic_irreducible_over}
\formallabel{GlobalTheory.exists_prime_dvd_sub_one_irreducible_cyclotomic}
\formallabel{GlobalTheory.irreducible_cyclotomic_of_not_dvd_discr}
\end{lemma}

\begin{proof}
By Dirichlet's theorem there is a prime $m\equiv1\pmod n$ with $m>|d_K|$, so $m\nmid d_K$. Let $A=\Q(\zeta_m)$ and let $F_0$ be
the intersection of $A$ and $K$ in a common extension. The discriminant of $F_0$ divides those of $A$ and of $K$; since
$d_A=\pm m^{m-2}$, $d_{F_0}$ divides a power of $m$ and is prime to $m$, so $|d_{F_0}|=1$ and $F_0=\Q$ by Minkowski's
theorem \cite[Ch.~III, Thm.~2.17]{Neukirch1999}. As $A/\Q$ is Galois, $A$ and $K$ are linearly disjoint, so $[K(\zeta_m):K]=
[A:\Q]=m-1=\deg\Phi_m$, and $\Phi_m$ is the minimal polynomial of $\zeta_m$ over $K$.
\end{proof}

\begin{lemma}[Crossing field]\label{pf:sh-hodge:cross}
Let $F,\Omega,L$ be as in \cref{pf:sh-hodge:cheb}, $g\in G_F$, and $m$ a prime with $[L:F]\mid m-1$ and $\Phi_m$ irreducible over
$L$. There are a finite extension $Z\subseteq\Omega$ of $F$ and $\tau\in(\Z/m)^\times$ such that $\Phi_m$ is irreducible over
$Z$ and every $\sigma\in G_Z$ with $\sigma(\zeta)=\zeta^\tau$ for all $\zeta\in\mu_m$ satisfies $\sigma g^{-1}\in G_L$.
\formallabel{GlobalTheory.T.exists_cyclotomic_crossing_field_deciding_frobenius}
\formallabel{GlobalTheory.ChebCross.cross}
\end{lemma}

\begin{proof}
Let $\zeta$ be a primitive $m$-th root of unity, $\chi\colon G_F\to(\Z/m)^\times$ the character with $\sigma\zeta=\zeta^{\chi(\sigma)}$,
and $\tau$ a generator of the cyclic group $(\Z/m)^\times$. Since $\Phi_m$ is irreducible over $L$, $\chi(G_L)=(\Z/m)^\times$, so
there is $u\in G_L$ with $\chi(u)=\tau\chi(g)^{-1}$; put $h=ug$, so that $h|_L=g|_L$ and $\chi(h)=\tau$. Let $M=L(\zeta)$, a finite
Galois extension of $F$, and $Z=M^{\langle h\rangle}$ the fixed field of $h$ in $M$. The powers $h^j$ fix $Z$ and $\chi(h^j)=\tau^j$
runs through $(\Z/m)^\times$, so $\Phi_m$ is irreducible over $Z$. Let $\sigma\in G_Z$ act on $\mu_m$ by $\tau$. Since
$\Gal(M/Z)$ is generated by $h|_M$, $\sigma|_M=h^k|_M$ for some $k$; comparing the actions on $\zeta$ gives $\tau^k=\tau$, so
$k\equiv1$ modulo the order $m-1$ of $\tau$. The order of $h|_L$ divides $[L:F]$, which divides $m-1$, hence $h^k|_L=h|_L$ and
$\sigma|_L=h|_L=g|_L$, i.e.\ $\sigma g^{-1}\in G_L$.
\end{proof}

\begin{lemma}[Primes of degree one in residue classes]\label{pf:sh-hodge:dirichlet}
Let $Z$ be a number field, $m$ a prime with $\Phi_m$ irreducible over $Z$, $\tau\in(\Z/m)^\times$ and $S$ a finite set of finite
places of $Z$. There is a place $w\notin S$ of $Z$ whose absolute norm $\shhodgeNm w$ is a prime number with $\shhodgeNm w\equiv\tau
\pmod m$.
\end{lemma}

This is \lref{GlobalTheory.T.exists_prime_absNorm_eq_of_irreducible_cyclotomic'} (about $1\,300$ lines, deferred). Its
formal proof is the classical one \cite{Hecke1920}: for a character $\chi\ne1$ of $(\Z/m)^\times$, irreducibility of $\Phi_m$
over $Z$ provides a prime $P$ of $Z$ with $\chi(\shhodgeNm P)\ne1$
(\lref{GlobalTheory.T.exists_prime_norm_cyclotomicChar_ne_one}), hence the $L$-function of the ideal character
$\chi\circ\shhodgeNm$ is continuous and non-zero at $s=1$
(\lref{GlobalTheory.T.exists_continuous_normCharLFunction_ne_zero_one}); together with the pole of the Dedekind zeta function
of $Z$ at $s=1$ (\lref{GlobalTheory.T.exists_entire_sub_one_mul_euler_exp_ne_zero_one}) this shows that $\sum1/\shhodgeNm P$
over the primes $P$ with $\shhodgeNm P\equiv\tau$ diverges (\lref{GlobalTheory.Mem.RayPrimes.not_summable_class}), whereas the
primes whose norm is not prime contribute a convergent sum.

\begin{proof}[Proof of \cref{pf:sh-hodge:cheb}]
Lift $g$ to $g'\in G_F$. By \cref{pf:sh-hodge:primem} (for $K=L$, $n=[L:F]$) choose a prime $m$ with $[L:F]\mid m-1$ and $\Phi_m$
irreducible over $L$, and let $Z,\tau$ be as in \cref{pf:sh-hodge:cross} for $g'$. The set $S'$ of places of $F$ which lie in $S$ or
have a place $B$ above them whose inertia group does not fix $L$ is finite, since only finitely many places ramify in $L$
(\lref{GlobalTheory.T.finite_places_inertia_not_fixing_finiteExtension}); so is the set $S''$ of places of $Z$ above $S'$. By
\cref{pf:sh-hodge:dirichlet} there is a place $w\notin S''$ of $Z$ with $p=\shhodgeNm w$ prime and $p\equiv\tau\pmod m$. Choose a place $B$
of $\Omega$ above $w$ and a Frobenius $\sigma\in G_Z$ at $B$ for $p$. Let $v=w\cap\OO_F$. Since $\shhodgeNm w=(\shhodgeNm v)^{f(w|v)}$ is
prime, $\shhodgeNm v=p$, $B$ lies above $v$, and $\sigma$, viewed in $G_F$, is a Frobenius at $B$ for $\shhodgeNm v$
(\lref{GlobalTheory.T.isFrob_restrictScalars_of_prime_absNorm'}); $v$ has prime norm, i.e.\ absolute residue degree $1$
(\lref{GlobalTheory.absNorm_prime_iff_inertiaDeg_eq_one_and_finite_ramified}). As $w\notin S''$, $v\notin S$ and
$\shhodgeIne_B$ fixes $L$.

It remains to show $\sigma|_L=g$. First, $m\notin w$: otherwise $\shhodgeNm w$ divides $m^{[Z:\Q]}$, so $p=m$, contradicting
$p\equiv\tau\not\equiv0$. Hence $m$ is a unit at $B$, and two $m$-th roots of unity $a\ne b$ are not congruent modulo the maximal ideal of
$B$, since $1-b/a$ divides $m=\prod_{\zeta'\in\mu_m\smallsetminus\{1\}}(1-\zeta')$. For $\zeta\in\mu_m$, $\sigma(\zeta)$ and
$\zeta^p=\zeta^\tau$ are $m$-th roots of unity, congruent since $\sigma$ is a Frobenius for $p$; so $\sigma(\zeta)=\zeta^\tau$. By the
choice of $Z$, $\sigma g'^{-1}\in G_L$, i.e.\ $\sigma|_L=g'|_L=g$.
\end{proof}

\subsubsection{Poitou--Tate duality from its normalised form}

Let $F$ be a number field with algebraic closure $\Omega$, $\ell$ a prime and $r\ge1$, and let $\chi_r\colon G_F\to(\Z/\ell^r)^\times$
be the cyclotomic character modulo $\ell^r$. For a place $A$ of $\Omega$ above a finite place of $F$ let $\mu_{A}$ be
$\Z/\ell^r$ with the action of $\shhodgeDec_A$ through $\chi_r$, and $Z^2(\shhodgeDec_A,\mu_A)$, $H^2(\shhodgeDec_A,\mu_A)$ the
continuous $2$-cocycles and the second cohomology. A \emph{family of invariant maps} is a family $\shhodgeinv=(\shhodgeinv_A)_A$
of $\Z/\ell^r$-linear maps $\shhodgeinv_A\colon H^2(\shhodgeDec_A,\mu_A)\to\Z/\ell^r$, one for every valuation ring $A$ of $\Omega$.

\begin{definition}[Normalised and compatible invariants]\label{pf:sh-hodge:def-inv}
Let $\zeta\in\Omega$ be a primitive $\ell^r$-th root of unity.
\begin{enumerate}[label=(\alph*)]
\item $\shhodgeinv_A$ is \emph{normalised} at $A$ above $v$ if it takes the value $1$ on the class of every cocycle
  $z(g,h)=\kappa(g)\chi(h)$ where: $\varpi\in\OO_F$ is a uniformiser at $v$ ($\varpi\in v\smallsetminus v^2$), $y\in\Omega$ satisfies
  $y^{\ell^r}=\varpi$, $\kappa\colon\shhodgeDec_A\to\Z/\ell^r$ is the Kummer map $g(y)=\zeta^{\kappa(g)}y$, and $\chi\colon
  \shhodgeDec_A\to\Z/\ell^r$ is a homomorphism vanishing on $\shhodgeIne_A$ and equal to $1$ on the Frobenius elements at $A$ for
  $\shhodgeNm v$.
\item The family is \emph{conjugation compatible} if for all field automorphisms $\sigma$ of $\Omega$ and $\tau$ of $F$ with
  $\sigma|_F=\tau$, all $c\in\Z/\ell^r$ with $\sigma(\zeta)=\zeta^c$, all places $A$ above some $v$ and $A'=\sigma(A)$, and all cocycles $z$
  on $\shhodgeDec_A$ and $z'$ on $\shhodgeDec_{A'}$ with $z'(g',h')=c\,z(g,h)$ whenever $g'\sigma=\sigma g$ and $h'\sigma=\sigma h$, we have
  $\shhodgeinv_{A'}[z']=\shhodgeinv_A[z]$.
\item The family is a family of \emph{local invariants} for $\zeta$ if every $\shhodgeinv_A$, for $A$ above some finite place $v$,
  is bijective and normalised, and the family is conjugation compatible.
\end{enumerate}
\formal{GlobalTheory.Published.PoitouTate.IsNormalised}
\formal{GlobalTheory.Published.PoitouTate.ConjCompatible}
\formal{GlobalTheory.Published.PoitouTate.IsLocalInvariants}
\end{definition}

The proposition $\sfname{PoitouTateDuality}$ of the development states: for all $F,\Omega,\ell,r\ge1$ with $\ell\ne2$ or $F$
totally complex, and every primitive $\ell^r$-th root of unity $\zeta$, there is a family of local invariants for $\zeta$ such that
for every finite set $S$ of places containing those above $\ell$, every choice of places $A_v$ above $v\in S$, every lift of $\chi_r$
to $G_{F,S}$, and every finite discrete continuous $\Z/\ell^r[G_{F,S}]$-module $M$, the nine-term Poitou--Tate sequence (built with
these invariants) is exact, the Tate--Shafarevich groups $\Sha^1$ and $\Sha^2$ of $M$ and of its dual are dual, and
$H^n(G_{F,S},M)=0$ for $n\ge3$ \cite[Ch.~I, Thm.~4.10]{Milne2006}, \cite[Ch.~VIII, \S6]{NeukirchSchmidtWingberg2008}. Its
\emph{normalised form} $\sfname{PoitouTateNormalised}$ states the same with ``family of local invariants'' replaced by ``family
whose members above finite places are bijective and normalised'', i.e.\ without conjugation compatibility.
\formal{GlobalTheory.Published.PoitouTateDuality}
\formal{GlobalTheory.Published.PoitouTateNormalised}

\begin{theorem}\label{pf:sh-hodge:pt}
$\sfname{PoitouTateNormalised}$ implies $\sfname{PoitouTateDuality}$. More precisely, every family of invariant maps whose members
above finite places are bijective and normalised is conjugation compatible.
\formallabel{GlobalTheory.PoitouTateConj.poitouTateDuality_of_normalised}
\formallabel{GlobalTheory.PoitouTateConj.conjCompatible}
\end{theorem}

We use two standard facts, stated in the development as published results: for a place $A$ above $v$ and $n\ge2$ there is a
continuous homomorphism $\chi\colon\shhodgeDec_A\to\Z/n$ vanishing on $\shhodgeIne_A$ and equal to $1$ on the Frobenius elements at $A$ for
$\shhodgeNm v$ (\lref{GlobalTheory.Published.PoitouTate.Model.exists_unramified}; $\shhodgeDec_A/\shhodgeIne_A$ is procyclic, topologically generated
by the Frobenius, \cite[Ch.~I, \S7]{Serre1979}); and for every continuous homomorphism
$\chi\colon\shhodgeDec_A\to\Z/\ell^r$ there are a uniformiser $\varpi$ at $v$, a root $y$ of $Y^{\ell^r}=\varpi$ and its Kummer map
$\kappa$, such that $z(g,h)=\kappa(g)\chi(h)$ is a $2$-cocycle on $\shhodgeDec_A$ with values in $\mu_A$
(\lref{GlobalTheory.Published.PoitouTate.Model.exists_kummer}; Kummer theory, \cite[Ch.~IV, \S3]{Neukirch1999}).

\begin{lemma}\label{pf:sh-hodge:placemap}
Let $\sigma,\tau$ be as in \cref{pf:sh-hodge:def-inv}(b), $A$ a place above $v$ and $A'=\sigma(A)$. Then $A'$ lies above the place
$v'=\tau(v)$, $\shhodgeNm v'=\shhodgeNm v$, and for every uniformiser $\varpi$ at $v$ the element $\varpi'=\tau(\varpi)$ is a uniformiser at $v'$
with $\varpi'=\sigma(\varpi)$ in $\Omega$. The map $\psi\colon\shhodgeDec_{A'}\to\shhodgeDec_A$, $g'\mapsto\sigma^{-1}g'\sigma$, is a continuous
homomorphism which maps $\shhodgeIne_{A'}$ into $\shhodgeIne_A$ and Frobenius elements at $A'$ for $q$ to Frobenius elements at $A$ for
$q$, and $\chi_r\circ\psi=\chi_r$.
\formallabel{GlobalTheory.PoitouTateConj.exists_place_map}
\formallabel{GlobalTheory.PoitouTateConj.conjBack_mem_Dec}
\formallabel{GlobalTheory.PoitouTateConj.conjBack_mem_Ine}
\formallabel{GlobalTheory.PoitouTateConj.conjBack_isFrob}
\formallabel{GlobalTheory.PoitouTateConj.cycMod_conjBack}
\formallabel{GlobalTheory.PoitouTateConj.continuous_conjBack}
\end{lemma}

\begin{proof}
$\tau$ restricts to a ring automorphism of $\OO_F$ mapping $v$ to the prime $v'=\tau(v)$, with $\OO_F/v\cong\OO_F/v'$; for
$x\in\OO_F$, $x\in v'$ iff $\tau^{-1}(x)\in v$ iff $\tau^{-1}(x)$ lies in the maximal ideal of $A$ iff $x=\sigma(\tau^{-1}x)$ lies in
that of $A'$. Uniformisers are mapped to uniformisers since $\tau$ maps $v^2$ to $v'^2$. For $g'\in\shhodgeDec_{A'}$, $\psi(g')$ fixes
$F$ and preserves $A$; the conditions defining inertia and Frobenius elements are transported by $\sigma$. The map $\psi$ is continuous for the
Krull topologies, since the preimage of the subgroup fixing a finite extension $M$ of $F$ contains the subgroup fixing $\sigma(M)$. Finally
$\psi(g')$ acts on $\sigma^{-1}(\mu_{\ell^r})=\mu_{\ell^r}$ by $\sigma^{-1}\circ(x\mapsto x^{\chi_r(g')})\circ\sigma$, which is again
$x\mapsto x^{\chi_r(g')}$.
\end{proof}

\begin{proof}[Proof of \cref{pf:sh-hodge:pt}]
Let $\shhodgeinv$ be a family whose members above finite places are bijective and normalised, and let
$\sigma,\tau,c,A,A',v,z,z'$ be as in \cref{pf:sh-hodge:def-inv}(b). Let $\psi$ be as in \cref{pf:sh-hodge:placemap}; then
$z'(g',h')=c\,z(\psi g',\psi h')$. For a cochain $w$ on $\shhodgeDec_A$ put $Tw(g',h')=c\,w(\psi g',\psi h')$. Since $\psi$ is a
continuous homomorphism preserving $\chi_r$, $T$ maps $2$-cocycles to $2$-cocycles, and $T(df)=d(c\,f\circ\psi)$ for every continuous
$1$-cochain $f$, so $T$ maps coboundaries to coboundaries.

By the first standard fact (with $n=\ell^r\ge2$) choose $\chi$ unramified with value $1$ on Frobenius elements, and by the second
choose $\varpi,y,\kappa$ and the cocycle $z_0=\kappa\cup\chi$. By normalisation $\shhodgeinv_A[z_0]=1$. Put
$a=\shhodgeinv_A[z]$; then $\shhodgeinv_A[z-az_0]=0$, so $z-az_0=df$ is a coboundary by injectivity. Hence
$z'-aTz_0=T(z-az_0)$ is a coboundary on $\shhodgeDec_{A'}$, and $\shhodgeinv_{A'}[z']=a\,\shhodgeinv_{A'}[Tz_0]$.

We show that $Tz_0$ is a normalising cocycle at $A'$, so that $\shhodgeinv_{A'}[Tz_0]=1$. With $v',\varpi'$ as in
\cref{pf:sh-hodge:placemap}, $y'=\sigma(y)$ satisfies $y'^{\ell^r}=\sigma(\varpi)=\varpi'$; for $g'\in\shhodgeDec_{A'}$,
\[
g'(y')=\sigma(\psi(g')y)=\sigma(\zeta^{\kappa(\psi g')}y)=\zeta^{c\,\kappa(\psi g')}y',
\]
so $\kappa'=c\,\kappa\circ\psi$ is the Kummer map of $y'$; $\chi'=\chi\circ\psi$ is a homomorphism which vanishes on $\shhodgeIne_{A'}$
and equals $1$ on the Frobenius elements at $A'$ for $\shhodgeNm v'=\shhodgeNm v$, by \cref{pf:sh-hodge:placemap}; and
$Tz_0(g',h')=c\,\kappa(\psi g')\chi(\psi h')=\kappa'(g')\chi'(h')$. Hence $\shhodgeinv_{A'}[z']=a=\shhodgeinv_A[z]$, which is conjugation
compatibility.

Given the family provided by $\sfname{PoitouTateNormalised}$ for $(F,\Omega,\ell,r,\zeta)$, it is therefore a family of local invariants,
and the remaining assertions (exactness of the nine-term sequence, duality of $\Sha$, vanishing of $H^n$ for $n\ge3$) are the same
in both propositions.
\end{proof}

% ---------------------------------------------------------------------------
\subsection{Local Galois cohomology at all places}\label{pf:sh-hodge:bk}

Throughout this subsection $\mathrm{PD}$ is a period datum for $G=G_K$, $K/\Ql$ finite, and $\chi$ its cyclotomic character. For a
good coefficient ring $(\mathit{LL},\OO)$ we write $\mathit{LL}(1)$ for $\mathit{LL}$ with $G$ acting through $\chi$ and $\OO(1)\subseteq
\mathit{LL}(1)$ for the lattice $\OO$. A bilinear map $P\colon V\times V'\to\mathit{LL}$ between representations is a \emph{twisted duality} if
$P(gx,gy)=\chi(g)P(x,y)$, $P$ is non-degenerate on both sides and $\dim V=\dim V'$; lattices $T\subseteq V$, $T'\subseteq V'$ are
\emph{dual} for $P$ if moreover $T'=\{y:\ P(T,y)\subseteq\OO\}$. For cochains $a\colon G\to V$, $b\colon G\to V'$ put $(a\cup b)(g,h)=
P(a(g),g\,b(h))$; for $d\colon G\to T$ and $c\colon G\to V'/T'$ put $(d\cup c)(g,h)=P(d(g),g\,\tilde c(h))$ for any lift $\tilde c$, a
cochain with values in $\mathit{LL}(1)/\OO(1)$. A $2$-cochain with values in $\mathit{LL}(1)$ is a \emph{coboundary} if it is $df$ for a
continuous $f\colon G\to\mathit{LL}(1)$; a cochain with values in $\mathit{LL}(1)/\OO(1)$ is a coboundary if it is $df$ for some
$f\colon G\to\mathit{LL}(1)/\OO(1)$.

\begin{definition}[The Bloch--Kato statements]\label{pf:sh-hodge:def-bk}
The following statements about $\mathrm{PD}$ quantify over all good coefficient rings $(\mathit{LL},\OO)$, all representations $V,V'$
over $\mathit{LL}$ with $V$ de~Rham, and all twisted dualities $P\colon V\times V'\to\mathit{LL}$.
\begin{enumerate}
\item[(A)] \emph{Rational duality.} A continuous cocycle $a$ of $V$ is finite if and only if $a\cup b$ is a coboundary for every finite
  cocycle $b$ of $V'$; and symmetrically for cocycles of $V'$.
\item[(Ad)] \emph{Dimension formula.} $\dim\shhodgeHf(V)=\dim H^0(G,V)+\dim D_{\mathrm{dR}}(V)-\dim\mathrm{Fil}^0D_{\mathrm{dR}}(V)$.
\item[(RE)] \emph{Right exactness.} If $f\colon V\to V'$ is a surjective $G$-map between crystalline representations, every finite
  cocycle of $V'$ is cohomologous to $f\circ c$ for a finite cocycle $c$ of $V$.
\item[(B)] \emph{Integral duality.} For dual lattices $T,T'$: a locally constant cocycle $c$ of $V'/T'$ lies in $\shhodgeHfD(T')$ if and
  only if $d\cup c$ is a coboundary for every $d\in\shhodgeZf(T)$; and $d\in Z^1(T)$ lies in $\shhodgeZf(T)$ if and only if $d\cup c$ is a
  coboundary for every locally constant $c\in\shhodgeHfD(T')$.
\item[(C)] \emph{Finite level duality.} For dual lattices $T,T'$, a uniformiser $\varpi$ of $\OO$ and $r\ge1$: a cocycle $c$ of
  $T/\varpi^rT$ is finite if and only if $\varpi^{-r}P(\tilde c,\tilde c')$ (computed on lifts) is a coboundary modulo $\OO(1)$ for every
  finite cocycle $c'$ of $T'/\varpi^rT'$; and symmetrically.
\end{enumerate}
$\mathrm{bk}(\mathrm{PD})$ denotes the conjunction of the five statements. The hypothesis \shhodgeBK{} asserts, for every tied local
theory and every $K$: (Ad) for $\Ql$-representations (with $\mathit{LL}=\Ql$), (A) for $\Ql$-representations, and (B) for $\Zl$-lattices
in $\Ql$-representations, with $\Zl(1)\subseteq\Ql(1)$ in place of $\OO(1)$.
\formal{Formalism.MainNotions.LocalSituation.bkAll}
\formal{Formalism.MainNotions.N.bkACore}
\formal{Formalism.MainNotions.N.bkBCore}
\formal{GlobalTheory.Published.BlochKato.DimFormula}
\formal{GlobalTheory.Published.BlochKato.RationalDuality}
\formal{GlobalTheory.Published.BlochKato.IntegralDuality}
\end{definition}

In \cite{BlochKato1990} the statements over $\Ql$ are Cor.~3.8.4 (dimension formula) and Prop.~3.8 (the subspaces $H^1_f$ of $V$
and $V^*(1)$ are exact annihilators of each other under local Tate duality); the integral statement is the corresponding
statement for $H^1_f$ of lattices and of $V/T$ (compare \cite[\S3]{BlochKato1990}). The point of this subsection is to pass from $\Ql$
to arbitrary coefficients and to finite level.

\begin{theorem}[Bloch--Kato at all places]\label{pf:sh-hodge:bkall}
Assume \shhodgeBK, \shhodgeLTD{} and \shhodgeLEC. For every datum $D$, every prime $\ell$ and every finite extension $K/\Ql$, the
statements $\mathrm{bk}(\mathrm{PD}^D_K)$ hold.
\formallabel{GlobalTheory.BlochKatoLocalOfHyp.blochKatoLocal_allPrimes_of_blochKato}
\formallabel{GlobalTheory.T.blochKatoLocal_all_of_tateDuality_eulerChar}
\formallabel{GlobalTheory.BlochKatoLocalOfHyp.bkAll_of_blochKatoLocal}
\end{theorem}

\begin{remark}\label{pf:sh-hodge:bkrem}
The two formal statements labelled here are identical (the second is proved by quoting the first); the conclusion is called
\emph{Conclusion~5} in the development. Only the tied local theory of $D$ at $\ell$ enters; \shhodgeLTD{} and \shhodgeLEC{} are used only
for (C) (and \shhodgeLEC{} not even there, see below).
\end{remark}

The proof treats the five statements separately. Fix $D,\ell,K$, write $\mathrm{PD}=\mathrm{PD}^D_K$, and let $(\mathit{LL},\OO)$ be a good
coefficient ring. Every representation $V$ over $\mathit{LL}$ is also a $\Ql$-representation $V_{\Q}$ (restriction of scalars); a cocycle
of $V$ is finite iff it is finite as a cocycle of $V_\Q$ (the condition $c(g)\otimes1=gy-y$ in $V\otimes_{\Ql}B$ does not involve
$\mathit{LL}$), $V$ is de~Rham iff $V_\Q$ is, and $D_{\mathrm{dR}}(V_\Q)$, $\mathrm{Fil}^0D_{\mathrm{dR}}(V_\Q)$, $H^0(G,V_\Q)$ and
$\shhodgeHf(V_\Q)$ are the corresponding $\mathit{LL}$-spaces for $V$, viewed over $\Ql$
(\lref{GlobalTheory.PstCoeff.isDR_resQ_iff}).

\begin{lemma}[(Ad) with coefficients]\label{pf:sh-hodge:bkdim}
The $\Ql$-dimension formula of \shhodgeBK{} implies (Ad).
\formallabel{GlobalTheory.T.blochKatoDim_of_Ql_dimensionFormula}
\formallabel{GlobalTheory.bkADim_of_dimFormula}
\end{lemma}

\begin{proof}
Apply the formula to $V_\Q$. By the remarks above each of the four $\Ql$-dimensions is $[\mathit{LL}:\Ql]$ times the corresponding
$\mathit{LL}$-dimension (for $\shhodgeHf$ through the identification of finite cocycles and coboundaries); divide by
$[\mathit{LL}:\Ql]\ne0$.
\end{proof}

\begin{lemma}[(A) with coefficients]\label{pf:sh-hodge:bkA}
The rational duality of \shhodgeBK{} implies (A).
\formallabel{GlobalTheory.T.blochKatoAnnihilator_of_rationalLocalDuality}
\formallabel{GlobalTheory.bkACore_of_rationalDuality}
\end{lemma}

\begin{proof}
Let $P$ be a twisted duality over $\mathit{LL}$ and $P_\Q=\shhodgeTr_{\mathit{LL}/\Ql}\circ P$. Then $P_\Q$ is a twisted duality between $V_\Q$
and $V'_\Q$, since the trace form of the separable extension $\mathit{LL}/\Ql$ is non-degenerate
(\lref{GlobalTheory.BkCoeff.isDual1_trPair}), and $a\cup_{P_\Q}(\lambda b)=\shhodgeTr(\lambda\,(a\cup_Pb))$ for $\lambda\in\mathit{LL}$. An
$\mathit{LL}(1)$-valued $2$-cochain $z$ is a coboundary iff $\shhodgeTr(\lambda z)$ is a $\Ql(1)$-coboundary for every $\lambda\in\mathit{LL}$
(\lref{GlobalTheory.BkCoeff.cobL_iff_trace}; write $z$ in a $\Ql$-basis of $\mathit{LL}$ and use the dual basis for the trace form).
Let $a$ be a cocycle of $V$. By the $\Ql$-statement, $a$ is finite iff $a\cup_{P_\Q}b'$ is a coboundary for all finite $b'$ of $V'_\Q$. If
this holds and $b$ is a finite cocycle of $V'$, then $\lambda b$ is finite for every $\lambda$, so all $\shhodgeTr(\lambda(a\cup_Pb))$ are
coboundaries and $a\cup_Pb$ is a coboundary. Conversely, if $a\cup_Pb$ is a coboundary for all finite $b$, then so is
$a\cup_{P_\Q}b=\shhodgeTr(a\cup_Pb)$, and $a$ is finite. The symmetric statement is proved in the same way.
\end{proof}

\begin{proof}[Proof of \cref{pf:sh-hodge:bkall}]
By \shhodgeBK{} applied to the tied local theory of $D$ at $\ell$ we have the three $\Ql$-statements at $K$.
(Ad) and (A) are \cref{pf:sh-hodge:bkdim,pf:sh-hodge:bkA}. (RE) holds for every tied local theory
(\lref{GlobalTheory.BKRightExact.bkARightExact_of_tied}, whose proof lifts the periods of a finite cocycle of $V'$ along $f$, using
the tie with Fontaine's rings and the fundamental exact sequence $0\to\Ql\to B_{\mathrm{cris}}^{\varphi=1}\oplus\mathrm{Fil}^0B_{\mathrm{dR}}
\to B_{\mathrm{dR}}$). (B) follows from the integral statement over $\Zl$ by
\cref{pf:sh-hodge:bkB}, and (C) from (A), (B) and \shhodgeLTD{} by \cref{pf:sh-hodge:bkC}.
\end{proof}

\begin{proposition}[(B) with coefficients]\label{pf:sh-hodge:bkB}
The integral duality of \shhodgeBK{} over $\Zl$ implies (B).
\formallabel{GlobalTheory.T.blochKato_annihilator_of_integralLocalDuality}
\formallabel{GlobalTheory.BKIntDescent.bkBCore_of_integralDuality}
\end{proposition}

\begin{proof}
The argument is a descent from $\OO$ to $\Zl$ through a dual functional. There are a $\Ql$-linear $\lambda\colon\mathit{LL}\to\Ql$ with
$\lambda(\OO)\subseteq\Zl$ and elements $e_1,\dots,e_n,f_1,\dots,f_n\in\OO$ with $x=\sum_i\lambda(f_ix)e_i$ for all $x\in\mathit{LL}$
(\lref{GlobalTheory.T.exists_integral_dual_functional}, deferred); in particular $x\in\OO$ iff $\lambda(ox)\in\Zl$ for all $o\in\OO$. Regard $T,T'$ as $\Zl$-lattices of $V_\Q,V'_\Q$
(\lref{GlobalTheory.T.zlLattice_of_coefficientLattice}, deferred) and put $P_\Z=\lambda\circ P$. Then $P_\Z$ is a twisted duality and $T,T'$
are dual for $P_\Z$ over $\Zl$, by the choice of $\lambda$ (\lref{GlobalTheory.BKIntDescent.dualLatZ}, deferred). Moreover a cochain $z$
with values in $\mathit{LL}(1)/\OO(1)$ is a coboundary iff $\lambda(oz)$ is a coboundary in $\Ql(1)/\Zl(1)$ for every $o\in\OO$
(\lref{GlobalTheory.BKIntDescent.cobO_iff}, deferred), and $d\cup_{P_\Z}c$ is $\lambda(d\cup_Pc)$ up to a $\Zl(1)$-valued cochain. With these
dictionary entries the argument is that of \cref{pf:sh-hodge:bkA}: for the first assertion, if $c$ is finite then $\lambda(o\,(d\cup c))=
\lambda((od)\cup c)$ is a coboundary for every finite $d$ and every $o$, because $od$ is finite and by the $\Zl$-statement; conversely if $d\cup c$
is a coboundary for every finite $d$, so is $\lambda(d\cup c)=d\cup_{P_\Z}c$, and $c$ is finite by the $\Zl$-statement. The second assertion is
proved in the same way, using that $\shhodgeHfD(T')$ is an $\OO$-module.
\end{proof}

\begin{proposition}[(C), finite level]\label{pf:sh-hodge:bkC}
Assume \shhodgeLTD. If (A) and (B) hold for $\mathrm{PD}$, so does (C).
\formallabel{GlobalTheory.T.breuilKisin_partC_of_partA_partB_finite_level}
\formallabel{GlobalTheory.BKFinLevel.bkC_of_bkA_bkB}
\end{proposition}

\begin{proof}
Fix dual lattices $T,T'$, $\varpi$ and $r$. Write $\iota\colon T/\varpi^rT\to V/T$ for $x\mapsto\varpi^{-r}x$ (the \emph{divisible image}),
and similarly for $T'$. We use five lemmas of the development, whose proofs are deferred: (i) finite cocycles of $T/\varpi^rT$ and of
$T'/\varpi^rT'$ are orthogonal for the pairing of (C), a consequence of (A)
(\lref{GlobalTheory.T.blochKato_reductions_orthogonal_mod_pow}); (ii) a cocycle of $T'/\varpi^rT'$ (or of $T/\varpi^rT$) whose divisible
image lies in $\shhodgeHfD(T')$ (resp.\ $\shhodgeHfD(T)$) is finite (\lref{GlobalTheory.T.blochKato_quotient_of_divisible_image_mem}); (iii) for $d\in Z^1(T)$ and a cocycle $c'$ of
$T'/\varpi^rT'$, $d\cup\iota(c')$ is a coboundary iff the finite-level pairing of the reduction of $d$ with $c'$ is
(\lref{GlobalTheory.T.cup_divisible_image_coboundary_iff_finite_level}), and similarly with the roles exchanged
(\lref{GlobalTheory.T.finite_level_cup_swap_coboundary_iff}); (iv) \emph{Tate's double annihilator}: by \shhodgeLTD, if $S$ is an
$\OO$-submodule of $Z^1(V/T)$ containing $B^1(V/T)$ and $x\in Z^1(V/T)$ is orthogonal to every $d\in Z^1(T')$ orthogonal to $S$, then
$x\in S$ (\lref{GlobalTheory.T.lattice_dual_cyclotomic}).

Orthogonality is (i). Let $c'$ be orthogonal to all finite $c$. For $d\in\shhodgeZf(T)$ the reduction of $d$ is finite, so by (iii) $d\cup
\iota(c')$ is a coboundary; by (B), $\iota(c')\in\shhodgeHfD(T')$, and $c'$ is finite by (ii). Let now $c$ be orthogonal to all finite $c'$, and
$S=\iota(\text{finite cocycles of }T/\varpi^rT)+B^1(V/T)$. If $d\in Z^1(T')$ is orthogonal to $S$, then by (iii) its reduction modulo
$\varpi^r$, a cocycle of $T'/\varpi^rT'$, is orthogonal to all finite cocycles of $T/\varpi^rT$, hence finite by the previous step, and then
(iii) and the hypothesis on $c$ show that $\iota(c)$ is orthogonal to $d$ (for the transposed pairing, for which $T',T$ are dual). By (iv), $\iota(c)\in S$: $\iota(c)=\iota(a)+b$ with $a$ finite and $b\in B^1(V/T)$. Then $\iota(c-a)\in\shhodgeHfD(T)$, so $c-a$ is finite by
(ii) (for $T$), and $c$ is finite.
\end{proof}

\begin{remark}\label{pf:sh-hodge:bkC-rem}
The hypothesis \shhodgeLEC{} of the formal statement of \cref{pf:sh-hodge:bkC} is not used in its proof, nor is the tie of the local
theory, which is a hypothesis of the formal statement.
\end{remark}

\begin{corollary}[Bloch--Kato annihilators at all rows]\label{pf:sh-hodge:bka}
Assume \shhodgeBK, \shhodgeLTD{} and \shhodgeLEC. For every datum $D$, every local index $\iota$, every place $w$ of $E$ above
$\ell_\iota$ and every coefficient extension $\mathit{LL}$ at $w$ (a finite extension of $\Q_{\ell_\iota}$ containing the coefficient field
of $D$ at $w$), statement (A) holds for $\mathrm{PD}^D_{K_w}$ with the coefficient ring $(\mathit{LL},\OO_{\mathit{LL}})$.
\formallabel{GlobalTheory.T.blochKatoAnnihilator_all_of_tateDuality_eulerChar}
\formallabel{GlobalTheory.BKaRows.bkaAll_of_hypotheses}
\formallabel{Formalism.MainNotions.N.bloch_kato_annihilator_above_l}
\end{corollary}

\begin{proof}
By \cref{pf:sh-hodge:bkall}, (A) holds for $\mathrm{PD}^D_{K_w}$ and every good coefficient ring; specialise it to $(\mathit{LL},\OO_{\mathit{LL}})$, which
is good (\lref{GlobalTheory.T.padicIntegers_finiteExt_goodCoefficientRing}).
\end{proof}

\begin{remark}
The formal statement of \cref{pf:sh-hodge:bka} quantifies over \emph{good} coefficient extensions (containing the roots of unity of
order dividing $\lcm(\ell,N)$, \cref{sec:pf:linkN-adm}); this hypothesis is not used.
\end{remark}

\begin{proposition}[Finiteness of local cohomology]\label{pf:sh-hodge:fin}
Assume \shhodgeLTD. For every datum $D$, every local index $\iota$, every place $w$ of $E$ above $\ell_\iota$, every coefficient
extension $\mathit{LL}$ at $w$, every representation $V$ of $G_{K_w}$ over $\mathit{LL}$ and every lattice $T$ of $V$ over
$\OO_{\mathit{LL}}$, the lattice cohomology $H^1(G_{K_w},T)$ is a finitely generated $\OO_{\mathit{LL}}$-module.
\formallabel{GlobalTheory.T.tateFinitenessAll_of_localTateDuality}
\formallabel{GlobalTheory.finAll_of_localTateDuality}
\end{proposition}

\begin{proof}
By \lref{GlobalTheory.T.finite_H1_period_datum}, \shhodgeLTD{} implies that for every local theory at every prime and every finite
$K$, $H^1$ of every lattice over every good coefficient ring is finitely generated (compare \cite[Ch.~II, \S5.2]{Serre1997}). Apply
this to the local theory of $D$ at
$\ell_\iota$, $K=K_w$ and the good coefficient ring $(\mathit{LL},\OO_{\mathit{LL}})$
(\lref{GlobalTheory.T.padicIntegers_finiteExt_goodCoefficientRing}).
\end{proof}

% ---------------------------------------------------------------------------
\subsection{Hodge--Tate, Fontaine and comparison data}\label{pf:sh-hodge:hw}

\paragraph{Good triples and Wach structures.}
Fix a local theory at $\ell$ with $A_{\inf}$-data. A triple $(K,s,\mathit{LL})$ of a finite extension $K/\Ql$, an integer $s$ and a finite
extension $\mathit{LL}/\Ql$ is \emph{good} if $s\ge1$, $s\ge2$ when $\ell=2$, $K\subseteq K_s=K_0(\mu_{\ell^s})$ and $\mathit{LL}$ contains all
conjugates of $K_s$; it is in \emph{case~A} if moreover $K$ is unramified, $s=1$ and $\ell\ne2$. A \emph{Wach structure} for $(K,s,\mathit{LL})$ is the
axiomatised $(\varphi,\Gamma)$-module theory of \cref{pf:linkN-adm:wach}: the ring $A=\OO_{\mathit{LL}}[[\pi]]^{\Z/n}$ ($n=[K_0:\Ql]$ letters) with
$\varphi$ and a quotient $\Gamma$ of $G_K$ (with elements $\gamma$, $\gamma_K$ and the torsion subgroup $\Delta$), the ring $\mathcal E\supseteq A$,
for every lattice $T$ the \'etale module $D(T)$ and the Wach module $N(T)\subseteq D(T)$, the class maps $\shhodgecl$ (from Herr cocycles for
$\gamma$ to $H^1(G_{K_s},T)$) and $\shhodgecl_K$ (from $\Delta$-invariant Herr cocycles for $\gamma_K$ to $H^1(G_K,T)$), the comparison map
$N(T)\to D_{\mathrm{cris}}(V)$ and the level maps at $\ell$-power roots of unity. We write $\pi$ for the variable of every letter,
$q_0=((1+\pi)^\ell-1)/\pi$, $X=\varphi^{s-1}(\pi)$, $Q=\varphi^{s-1}(q_0)$, and $N_i(T)$ for the part of $N(T)$ at the letter $i$; $\pi_0\in A$ is
the element of the structure used in \cref{sec:pf:linkN-adm} (with $q_0$ a unit multiple of $\pi_0+\ell$ when $K=K_0$, $s=1$, $\ell\ne2$, by (P9)
below).
\formal{Formalism.BasicNotions.LocWach}

\begin{definition}[Hodge--Tate, Fontaine and comparison data]\label{pf:sh-hodge:def-hw}
Let $\mathcal L=(\mathrm{PD}^{(\ell)})_\ell$ be local theories at all primes, with $A_{\inf}$-data $\mathcal A=(A^{(\ell)})_\ell$.
\begin{enumerate}[label=(\alph*)]
\item \emph{Hodge--Tate data} $HW$ assign to every prime $\ell$ and every good triple $(K,s,\mathit{LL})$ a Wach structure
  $W(K,s,\mathit{LL})$ for $\mathrm{PD}^{(\ell)}_K$ with coefficients $\OO_{\mathit{LL}}$, of level $s$, all of whose Wach modules are
  projective over $\Zl[[\pi]]$ (acting diagonally on the letters).
\item \emph{Fontaine data} for $HW$ assign to every good triple an \emph{\'etale tie} of $W(K,s,\mathit{LL})$: $\OO_{\mathit{LL}}$-algebra maps
  $A\to\OO_{\mathit{LL}}\otimes A_{\inf}$ and $\mathcal E\to R_W=\OO_{\mathit{LL}}\otimes W(\C_\ell^\flat)$ sending $\pi$ to
  $1\otimes([\eps^{1/\ell^s}]-1)$, and for every lattice $T$ an injective $\OO_{\mathit{LL}}$-linear map $j_T\colon D(T)\to R_W\otimes T$, semilinear
  for $\mathcal E\to R_W$, compatible with $\varphi$ and with $G_K$ (acting on $D(T)$ through $\Gamma$), natural in $T$, whose image spans
  $R_W\otimes T$; together with explicit cocycle formulas for $\shhodgecl$ and $\shhodgecl_K$ in terms of $j_T$, and the compatibility of
  $A\to\OO_{\mathit{LL}}\otimes A_{\inf}$ with Teichm\"uller representatives of roots of unity of order prime to $\ell$. This is Fontaine's
  description of \'etale $(\varphi,\Gamma)$-modules \cite{Fontaine1990}.
\item \emph{Comparison data} for $(HW,\text{Fontaine data})$ assign to every good triple a \emph{comparison tie}: maps from $N(T)$ to the
  $\tau$-components of $D_{\mathrm{cris}}(V)$ at every level and root of unity, compatible with the level maps, with $\varphi$, with $\Gamma$ and
  with the comparison map, and a map $\beta_{\mathrm{dR}}\colon(\OO_{\mathit{LL}}\otimes A_{\inf})\otimes T\to V\otimes B_{\mathrm{dR}}$
  relating them to $j_T$ modulo the filtration; and the two facts $\iota_A(\mu)-t\in\mathrm{Fil}^2$ for $\mu=[\eps]-1$, and
  $\iota_A([\eps^{1/\ell}]-1)$ invertible in $\mathrm{Fil}^0B_{\mathrm{dR}}$.
\item \emph{Ties} for these data are, for every prime $\ell$, a map $\theta\colon B_{\mathrm{dR}}\to\C_\ell$ which is a surjective ring
  homomorphism on $\mathrm{Fil}^0$ with kernel $\mathrm{Fil}^1$, Galois equivariant, the identity on $\Omega_\ell$, integral on $A_{\inf}$, compatible
  with Teichm\"uller lifts, with $\mathrm{Fil}^{a+1}=t\,\mathrm{Fil}^a$ and $\iota_A(\mu)-t\in\mathrm{Fil}^2$.
\end{enumerate}
\formal{Formalism.BasicNotions.HWOf}
\formal{Formalism.BasicNotions.FonOf}
\formal{Formalism.BasicNotions.CmpOf}
\formal{Formalism.MainNotions.LocalSituation.Ties}
\end{definition}

\begin{definition}[The theorem of Bhatt--Scholze]\label{pf:sh-hodge:def-bs}
$\sfname{BhattScholzeCrystallineLattices}$ is the following statement, for every prime $\ell$, every finite $K/\Ql$ and the
$A_{\inf}$-prism of $\OO_K$: (i) for every tied local theory with $A_{\inf}$-data and every crystalline lattice $T$ of a representation of
$G_K$ over a finite extension of $\Ql$, there is a prismatic $F$-crystal over $\OO_K$ \emph{realising} $T$, i.e.\ with an identification of
its value on $A_{\inf}$, after extension to $W(\C_\ell^\flat)$, with $W(\C_\ell^\flat)\otimes_{\Zl}T$ compatible with Frobenius and the Galois
actions; (ii) every $G_K$-equivariant map between realised $\Zl$-lattices comes from a unique map of $F$-crystals. This is a form of the
equivalence between prismatic $F$-crystals on $\OO_K$ and crystalline $\Zl$-lattices \cite[Thm.~1.2 and \S5]{BhattScholze2023}.
\formal{GlobalTheory.Published.BhattScholzeCrystallineLattices}
\end{definition}

\begin{theorem}[Existence of the data]\label{pf:sh-hodge:exhw}
Assume $\sfname{BhattScholzeCrystallineLattices}$. For all local theories $\mathcal L$ with $A_{\inf}$-data $\mathcal A$ which are tied at every
prime, there are Hodge--Tate data $HW$ and Fontaine data for $HW$ admitting comparison data and ties.
\formallabel{GlobalTheory.HWL.exists_hw_fon_cmp_ties}
\end{theorem}

The proof combines three constructions, made separately at each prime $\ell$; their proofs are deferred.
\begin{enumerate}[label=(\alph*)]
\item \emph{Wach lattices exist at all layers} (\lref{GlobalTheory.T.wachLatticesExistAllLayers_of_bhattScholze}, about $2\,350$ lines):
  under the theorem of Bhatt--Scholze (\cref{pf:sh-hodge:def-bs}), for every tied theory, every good triple and every choice of the auxiliary
  layers of a
  Wach structure (a $\Gamma$-structure, a ring of Wach series, an embedding layer into $A_{\inf}$ and an \'etale module), every crystalline
  lattice has a Wach module, i.e.\ an $A$-submodule of $D(T)$ with the properties of $N(T)$. The construction evaluates the $F$-crystal realising
  $T$ on a \emph{Wach prism} over $\OO_K$, a prism structure on a ring of power series in $\pi$
  (\lref{GlobalTheory.T.exists_wach_prism_over_integers}), and descends the result to $A$ (\cite{Wach1996,Berger2004} for the classical
  theory of Wach modules).
\item \emph{Faithfulness of the cyclotomic kernel} (\lref{GlobalTheory.T.cyclotomicKernel_tilt_invariants_faithful}): if $g\in G_{\Ql}$ has
  $\chi(g)=1$ and fixes every element of $\OO_{\C_\ell^\flat}$ fixed by $\{h\in G_L:\chi(h)=1\}$, then $g\in G_L$.
\item \emph{The local Wach functor} (\lref{GlobalTheory.T.exists_local_wach_functor}, about $5\,700$ lines): for a tied theory satisfying the
  conclusions of (a) and (b) there are Hodge--Tate data at $\ell$ and Fontaine data admitting comparison data. Its proof builds, for each good
  triple, a $\Gamma$-structure (\lref{GlobalTheory.HWL.GroupLayer.nonempty_gamS}), the ring of Wach series
  (\lref{GlobalTheory.HWL.Ser.nonempty_wSer}), the embedding into $A_{\inf}$ (\lref{GlobalTheory.HWL.IoLayer.main}), the \'etale modules
  from a Frobenius-stable basis of $R_W\otimes T$ fixed by the cyclotomic kernel (\lref{GlobalTheory.EtaleBasis.exists_etBasis},
  \lref{GlobalTheory.HWL.EtMod.nonempty_wEt}), the Wach modules from (a), the Herr class maps
  (\lref{GlobalTheory.HerrClassMaps.exists_classMaps}) and the comparison maps (\lref{GlobalTheory.HWL.Cmp.nonempty_cmpS}).
\end{enumerate}

\begin{lemma}\label{pf:sh-hodge:exhw-lem}
Let $\mathcal L,\mathcal A$ be tied local theories with $A_{\inf}$-data such that at every prime the conclusion of (a) holds. Then there are Hodge--Tate data,
Fontaine data and comparison data for $(\mathcal L,\mathcal A)$.
\formallabel{GlobalTheory.T.exists_hwOf_fonOf_cmpOf'}
\formallabel{GlobalTheory.HWL.exists_hwOf_fonOf_cmpOf}
\end{lemma}

\begin{proof}
At every prime $\ell$, (b) holds for every local theory, so (c) applies and gives Hodge--Tate data $HW_\ell$, Fontaine data $F_\ell$ and comparison data at
$\ell$. Choosing these for every $\ell$ gives the three families.
\end{proof}

\begin{proof}[Proof of \cref{pf:sh-hodge:exhw}]
By (a), the theory has Wach lattices at all layers at every prime, and \cref{pf:sh-hodge:exhw-lem} gives $HW$, Fontaine data and comparison data. The ties
are provided at every $\ell$ by the tie of the local theory with Fontaine's rings, through which Fontaine's map $\theta$ \cite{Fontaine1994} has the required
properties (\lref{GlobalTheory.T.nonempty_theta_deRhamFil_zero_concretePeriod}).
\end{proof}

\begin{theorem}[Herr's description of lattice cohomology]\label{pf:sh-hodge:herr}
For every datum $D$, all Hodge--Tate data $HW$ and Fontaine data $F$ of the local theory of $D$, every prime $\ell$ and every good triple
$(K,s,\mathit{LL})$, with Wach structure $W$:
\begin{enumerate}[label=(\alph*)]
\item if $\ell\ne2$ or $\sqrt{-1}\in K$: for every lattice $T$, $\shhodgecl_K$ is surjective onto $H^1(G_K,T)$, its kernel consists of the
  pairs $(\varphi m-m,\gamma_Km-m)$ with $m\in D(T)^\Delta$, and $\shhodgecl_K$ is natural for lattice maps;
\item for every lattice $T$, $\shhodgecl$ is surjective onto $H^1(G_{K_s},T)$ and its kernel consists of the pairs $(\varphi m-m,\gamma m-m)$,
  $m\in D(T)$;
\item $\shhodgecl$ is natural: for a lattice map $f$, $f_*\circ\shhodgecl=\shhodgecl\circ D(f)$.
\end{enumerate}
\formallabel{GlobalTheory.HerrFull.herr_part_a_every_record}
\end{theorem}

\begin{proof}
This is \lref{GlobalTheory.HerrFull.herr_a_and_natural}, which derives (a)--(c) from the statement of Fontaine and decompletion for the data
(\lref{GlobalTheory.HerrFull.fontaine_decompletion}) by Herr's computation of Galois cohomology through the complex
$D\xrightarrow{(\varphi-1,\gamma-1)}D\oplus D\to D$ \cite{Herr1998}.
\end{proof}

\paragraph{The later clauses of the Wach lattice theory.}
Let $W$ be a Wach structure of a good triple $(K,s,\mathit{LL})$, $p=\ell$, $u=\chi(\gamma)\in\OO^\times$, and let $T,T',T_0,T_1,T_2$ denote
crystalline lattices. For maps $M,h$ from the letters to $\N$, $T$ has jumps \emph{within $(M,h)$} if the Hodge--Tate jumps of $T$ at every
embedding over the letter $i$ lie in $[-M_i,h_i]$. The following clauses are the content of the norm bound for $\Delta$-invariant lattices
used in \cref{sec:pf:linkN-adm,sec:pf:linkN-sel}; each holds for all crystalline lattices.
\begin{enumerate}
\item[(P1)] \emph{Letters.} $N(T)=\sum_iN_i(T)$ and each $N_i(T)$ is free over $\OO[[\pi]]$ (on the $i$-th factor).
\item[(P2)] \emph{$\Gamma$-stability.} $N(T)$ is $\Gamma$-stable, $(g-1)N(T)\subseteq\pi N(T)$ for $g\in\Gamma$, and $(\gamma-1)N(T)\subseteq XN(T)$.
\item[(P3)] \emph{Comparison.} The comparison map vanishes only on elements of $\pi N(T)$, its image spans $D_{\mathrm{cris}}(V)$ over
  $\mathit{LL}$, and if $Q^k\varphi(x)\in N(T)$ its image is $p^k\varphi$ applied to the image of $x$.
\item[(P4)] \emph{Levels.} The level maps at the letters and at $p$-power roots of unity take values in the images of the $\tau$-components of
  $D_{\mathrm{cris}}$, vanish in negative degrees, span these images in degree $0$, and at the points of $HW$ are adapted to the Hodge
  filtration at $\tau$.
\item[(P5)] \emph{Heights.} If $T$ has jumps within $(M,h)$: $Q^{M_i}\varphi(N_{i-1}(T))\subseteq N(T)$, $Q^{h_i}N_i(T)\subseteq$ the span of
  $\varphi(N_{i-1}(T))$, and an element $x$ with $Q^kx\in N(T)$ for some $k$ lies in $N(T)$ iff $Q^kx$ lies in the vanishing ideals of the level
  maps at the roots of $Q$ times $N(T)$.
\item[(P6)] \emph{Saturation and exactness.} $N(T)$ is $p$-saturated in $D(T)$; a bijective lattice map induces an injection on Wach modules
  which is surjective up to powers of $p$; and a short exact sequence of crystalline lattices gives a sequence of Wach modules which is
  injective, exact and surjective up to powers of $p$, strict on the left, and has a cokernel of finite length.
\item[(P7)] \emph{Poles.} If $(y,z)$ is a Herr cocycle with $\prod_{j\in J}\varphi^j(Q)\,y\in N(T)$ for a multiset $J$ and $z\in\pi N(T)$, then
  $\shhodgecl(y,z)\in\shhodgeHf(G_{K_s},T)$.
\item[(P8)] \emph{Extensions of the unit.} For an extension $0\to T\to T_p\to\one\to0$ with $T_p$ crystalline and $T$ with jumps within $(M,h)$,
  and $x\in N(T_p)$ mapping to a multiple of an element of $N(\one)$ fixed by $\varphi$ and $\Gamma$: $\varphi x-x\in Q^{-\max M}N(T)$ and
  $\gamma x-x\in XN(T)$.
\item[(P9)] \emph{Tame descent.} If $s=1$ and $p\ne2$: $\#\Delta=(p-1)/[K:K_0]$; there are $\pi'\in A^\Delta$ and polynomials $E_i$ of degree
  $[K:K_0]$, coming from a polynomial over $K_0$ with a root in $K$ of valuation $1/[K:K_0]$, giving \emph{ring relations} ($A^\Delta$ is a
  power series ring in $\pi'$, $\pi'$ is a unit times $\pi^{(p-1)/[K:K_0]}$, and $q_0$, $\varphi(\pi')$ and $\gamma\pi'-\pi'$ are unit multiples of
  $E(\pi')$, $\pi'E(\pi')^{(p-1)/[K:K_0]}$ and $\pi'E(\pi')$); $N(T)$ has an $A$-basis in $N(T)^\Delta$ which spans $N(T)^\Delta$ over $A^\Delta$,
  and $(\gamma-1)N(T)^\Delta\subseteq\pi'N(T)^\Delta$; for $K=K_0$ one may take $\pi'=\pi_0$ with the polynomial $X+p$.
\item[(P10)] \emph{Characterisation.} $N(T)$ is the only free $A$-submodule $N_1$ of $D(T)$ with $\mathcal E\otimes N_1=D(T)$,
  $\gamma N_1\subseteq N_1$, $(\gamma-1)N_1\subseteq\pi N_1$, and such that for every $x\in N_1$ some $Q^k\varphi(x)$ lies in $N_1$ and some
  $Q^kx$ lies in the span of $\varphi(N_1)$.
\item[($\Gamma$ss)] For $y\ge1$ the quotient $Q^{-y}N(T)/N(T)$ is free of finite rank over $\OO$ and $\gamma$ is annihilated on it by
  $\prod_c(\gamma-c)$, $c$ running over the units $u^{-(a+j)}$ ($a$ a jump of $T$, $1\le j\le y$), each of which is an eigenvalue.
\item[(H1)--(H4)] \emph{Herr complex.} (H1) is \cref{pf:sh-hodge:herr}(b) and, if $s=1$ and $p\ne2$, (a); (H2) a class of $H^1(G_K,T)$ is
  finite iff its restriction to $G_{K_s}$ is; (H3) $\gamma(Q)=wQ$ with a unit $w\equiv1$ modulo $X$, and if $T$ has jumps within $(M,h)$, a Herr
  cocycle $(y,z)$ with $y\in Q^{-\max M}N(T)$ and $z\in XN(T)$ has a finite class, which vanishes only if $(y,z)$ is the coboundary of an element of
  $N(T)$; (H4) the analogous compatibility with the connecting map of a short exact sequence, up to a power of $p$.
\item[(L3)] \emph{Length lemma.} If $F\colon T_1\to T_2$ is onto and $T_1$ has jumps within $(M,h)$, then $X^{\max(M+h)}N_i(T_2)\subseteq
  F(N_i(T_1))$, and in case~$s=1$ a $\Delta$-version with $\pi'$.
\item[(E)] \emph{\'Etale comparison.} For the \'etale tie: $N(T)$ has an $A$-basis whose image under $j_T$ is an $R_W$-basis of $R_W\otimes T$,
  the $\varphi$-invariants of $R_W\otimes T$ are $1\otimes T$, and if $T$ has jumps within $(M,h)$ then $X^{\max h}$ maps $A_{\inf}\otimes T$ into
  $A_{\inf}N(T)$ and $X^{\max M}$ maps $A_{\inf}N(T)$ into $A_{\inf}\otimes T$.
\end{enumerate}
The precise forms are the formal definitions cited below; the descriptions above are complete except that (P4), (P9), (H4) and (L3) are
abbreviated.
\formal{Formalism.MainNotions.DeltaInvariantLattices.DeltaInvariantNormBound}
\formal{Formalism.MainNotions.DeltaInvariantLattices.UsedLocalConclusion}
\formal{Formalism.MainNotions.DeltaInvariantLattices.ConclusionE}

Four of the clauses have short proofs from other results, which we give.

\begin{proposition}[Clause (P1)]\label{pf:sh-hodge:P1}
For every datum $D$, all data $(HW,F)$, every good triple and every crystalline lattice $T$, $N(T)=\sum_iN_i(T)$ with $N_i(T)=e_iN(T)$ ($e_i$ the idempotent of the
letter $i$), and if $b_1,\dots,b_r$ is an $A$-basis of $N(T)$ then $e_ib_1,\dots,e_ib_r$ is a basis of $N_i(T)$ over $\OO[[\pi]]$ acting on the letter $i$. An
$A$-basis with $r=\rank_\OO T$ exists.
\formallabel{GlobalTheory.NL.LetterFree.letters_free}
\formallabel{GlobalTheory.NL.LetterFree.letters_loc}
\formallabel{GlobalTheory.NL.LetterFree.isLetterBasis}
\formallabel{GlobalTheory.NL.LetterFree.exists_letter_basis}
\formallabel{GlobalTheory.NL.LetterFree.lettersDecompose}
\end{proposition}

\begin{proof}
The existence of an $A$-basis of rank $\rank_\OO T$ is \lref{GlobalTheory.NL.N_basis_loc}, which deduces it from the finite generation of $N(T)$, the
freeness of $\mathcal E\otimes N(T)=D(T)$ of that rank, and the projectivity of $N(T)$ over $\Zl[[\pi]]$ required in the definition of Hodge--Tate data. Since
$\sum_ie_i=1$, $N(T)=\sum_ie_iN(T)$. If $y=\sum_ka_kb_k\in N(T)$ ($a_k\in A$) then $e_iy=\sum_k(a_k)_i\,e_ib_k$, where $(a_k)_i\in\OO[[\pi]]$ is the component at
$i$, since $e_ia_k$ is the element $(a_k)_i$ placed at the letter $i$. If $\sum_kc_k\,e_ib_k=\sum_kc'_k\,e_ib_k$ with $c_k,c'_k\in\OO[[\pi]]$, then
$\sum_k(c_k-c'_k)e_i\,b_k=0$ in $N(T)$, and the linear independence of the $b_k$ over $A$ gives $(c_k-c'_k)e_i=0$, i.e.\ $c_k=c'_k$.
\end{proof}

\begin{lemma}[Coboundaries on subgroups of finite index]\label{pf:sh-hodge:finind}
Let a group $G$ act linearly on a vector space $M$ over a field of characteristic $0$, $H\subseteq G$ a subgroup of finite index and $c\colon G\to M$ a
$1$-cocycle whose restriction to $H$ is a coboundary. Then $c$ is a coboundary.
\formallabel{GlobalTheory.NL.FiniteRestrictionH1f.coboundary_of_finiteIndex}
\formallabel{GlobalTheory.NL.FiniteRestrictionH1f.crisCob_of_finiteIndex}
\end{lemma}

\begin{proof}
Let $c(h)=hb-b$ for $h\in H$, and $d(g)=c(g)-(gb-b)$. Then $d$ is a cocycle vanishing on $H$, so $d(gh)=d(g)+g\,d(h)=d(g)$ and $d$ induces a function
$\bar d$ on $G/H$ with $\bar d(gq)=d(g)+g\,\bar d(q)$. Let $n=[G:H]$ and $S=\sum_{q\in G/H}\bar d(q)$. Summing over $q$ (and reindexing $q\mapsto gq$)
gives $S=n\,d(g)+gS$, so $d(g)=g(-S/n)-(-S/n)$ and $c(g)=gx-x$ with $x=b-S/n$.
\end{proof}

\begin{proposition}[Clause (H2)]\label{pf:sh-hodge:H2}
For every datum $D$, all data $(HW,F)$, every good triple and every lattice $T$ of a representation of $G_K$, a class $\kappa\in H^1(G_K,T)$ is finite if and
only if its restriction to $G_{K_s}$ is finite.
\formallabel{GlobalTheory.NL.FiniteRestrictionH1f.mem_H1f_iff_res_mem}
\formallabel{GlobalTheory.NL.FiniteRestrictionH1f.mem_H1fT_iff_res}
\formallabel{GlobalTheory.NL.FiniteRestrictionH1f.GKs_finiteIndex}
\end{proposition}

\begin{proof}
$G_{K_s}$ has finite index in $G_K$, since $[K_s:\Ql]=[K_0:\Ql](\ell-1)\ell^{s-1}$ is finite. If $\kappa$ is represented by $z$ with $z(g)\otimes1=gx-x$ in
$V\otimes B$, the restriction of $z$ satisfies the same identity on $G_{K_s}$. Conversely, let the restriction of $\kappa$ be represented by a finite cocycle
$c'$ of $G_{K_s}$, $c'(h)\otimes1=hx-x$, and let $z$ represent $\kappa$. Then $z(h)-c'(h)=hy-y$ for some $y\in T$ and all $h\in G_{K_s}$, so $g\mapsto z(g)
\otimes1$, a cocycle of $G_K$ in the $\Ql$-vector space $V\otimes B$, is a coboundary on $G_{K_s}$ (of $x+y\otimes1$). By \cref{pf:sh-hodge:finind} it is a
coboundary on $G_K$, i.e.\ $z$ is finite.
\end{proof}

\begin{proposition}[Clause (P6)]\label{pf:sh-hodge:P6}
For every datum $D$, all data $(HW,F,\mathrm{Cm})$ and every good triple, clause (P6) holds; in a short exact sequence $0\to T_0\xrightarrow{\iota}T_1
\xrightarrow{\beta}T_2\to0$ of crystalline lattices the sequence of Wach modules is even exact at $N(T_1)$.
\formallabel{GlobalTheory.T.wachModule_saturated_exactness_crystalline}
\formallabel{GlobalTheory.NL.Exact.exactness_of_N}
\formallabel{GlobalTheory.NL.Exact.length_cokN_lt_top}
\formallabel{GlobalTheory.NL.Exact.finite_of_pi_pow}
\formallabel{GlobalTheory.NL.Exact.finiteLength_of_killed}
\end{proposition}

\begin{proof}
Saturation of $N(T)$ and the statement on bijective lattice maps are \lref{GlobalTheory.NL.Isog.saturated_and_isogeny_invariant} (deferred). Let a short
exact sequence be given. The functor $D$ is exact on it (\lref{GlobalTheory.EtaleExact.DD_exact}, from the statement of Fontaine and decompletion), so
$D(\iota)$ is injective, which gives injectivity on $N(T_0)$; and $\ker D(\beta)=\operatorname{im}D(\iota)$. If $\iota(x)\in N(T_1)$ for some $x\in D(T_0)$ then
$x\in N(T_0)$ (\lref{GlobalTheory.T.mem_wachN_of_map_mem_saturated}, deferred); this is the strictness, and with $\ker D(\beta)=\operatorname{im}D(\iota)$ it gives
exactness at $N(T_1)$. There is $c$ such that $p^cz$ and $\pi^cz$ lie in $\beta(N(T_1))$ for every $z\in N(T_2)$
(\lref{GlobalTheory.T.wachN_cokernel_killed_by_power}, deferred, about $860$ lines); this is surjectivity up to $p^c$. The cokernel $C=N(T_2)/\beta(N(T_1))$ is
a finitely generated $A$-module killed by $\pi^c$; since $A/\pi^cA=\prod_{\text{letters}}\OO[[\pi]]/(\pi^c)$ is generated over $\OO$ by the elements $e_i\pi^k$
($e_i$ the idempotent of a letter, $k<c$), $C$ is finitely generated over $\OO$. It is killed by $p^c\ne0$, so it is a finitely generated module over
$\OO/p^c\OO$, an Artinian ring, and has finite length.
\end{proof}

\begin{proposition}[Clause (P8)]\label{pf:sh-hodge:P8}
For every datum $D$, all data $(HW,F,\mathrm{Cm})$ and every good triple, clause (P8) holds.
\formallabel{GlobalTheory.T.p8_of_cmpOf}
\formallabel{GlobalTheory.NL.NygaardFrobeniusPower.crystalline_extension_integral_cocycle_all_places}
\formallabel{GlobalTheory.NL.NygaardFrobeniusPower.core}
\formallabel{GlobalTheory.NL.NygaardFrobeniusPower.Q_pow_φ_mem}
\formallabel{GlobalTheory.NL.NygaardFrobeniusPower.isUnit_X}
\formallabel{GlobalTheory.NL.JumpsSES.jumps_ses_and_one}
\end{proposition}

\begin{proof}
Let $0\to T\xrightarrow{\iota}T_p\xrightarrow{\beta}\one\to0$ be exact with $T_p$ crystalline and $T$ with jumps within $(M,h)$, let $e\in N(\one)$ be fixed by
$\varphi$ and $\Gamma$, and let $x\in N(T_p)$ with $\beta(x)=ce$, $c\in\OO$. Sub- and quotient lattices of crystalline lattices are crystalline
(\lref{GlobalTheory.Per.CrysSubQuot.crys_sub_quot}), and in a short exact sequence of crystalline lattices the jumps of the middle term at $\tau$
are those of the two others (\lref{GlobalTheory.T.isJumpV_iff_of_exact_crystalline_lattices}); every jump of the unit object is $0\in[-M_i,h_i]$.
Hence $T_p$ has jumps within
$(M,h)$, and (P5) and (P2) apply to $T_p$ (\lref{GlobalTheory.T.wachLattice_heights_from_hodgeJumps},
\lref{GlobalTheory.NL.WachLatticeTrivialModUniformizer.wachLattice_galois_trivial_mod_uniformizer}). Writing $x=\sum_ie_ix$ with $e_ix\in N_i(T_p)$, the
first part of (P5) at the letter $i+1$ gives $Q^{M_{i+1}}\varphi(e_ix)\in N(T_p)$, hence $Q^{\max M}\varphi(x)\in N(T_p)$. We use that $D$ is exact (as in
\cref{pf:sh-hodge:P6}) and that $\iota(y)\in N(T_p)$ implies $y\in N(T)$.

Since $\beta(\varphi x-x)=c(\varphi e-e)=0$, there is $y\in D(T)$ with $\iota(y)=\varphi x-x$; then $\iota(Q^{\max M}y)=Q^{\max M}\varphi(x)-Q^{\max M}x\in N(T_p)$,
so $Q^{\max M}y\in N(T)$, i.e.\ $y\in Q^{-\max M}N(T)$. By (P2), $\gamma x-x=Xw$ with $w\in N(T_p)$. Then $X\beta(w)=\beta(\gamma x-x)=c(\gamma e-e)=0$; as $X=
\varphi^{s-1}(\pi)$ is a unit of $\mathcal E$ ($\pi$ is, and $\varphi$ is a ring endomorphism), $\beta(w)=0$, so $w=\iota(z)$ with $z\in D(T)$, and $z\in N(T)$ since
$\iota(z)=w\in N(T_p)$. Thus $\gamma x-x=\iota(Xz)$ with $Xz\in XN(T)$.
\end{proof}

\begin{theorem}[The later clauses of the Wach lattice theory]\label{pf:sh-hodge:wach}
For every datum $D$ and all Hodge--Tate, Fontaine and comparison data of its local theory, the clauses (P1)--(P10), ($\Gamma$ss), (H1)--(H4),
(L3) and (E) hold for the Wach structures of all good triples.
\formallabel{GlobalTheory.T.wachLattice_laterClauses_hold}
\formallabel{GlobalTheory.NL.wach_used_clauses}
\end{theorem}

\begin{proof}
The theorem is the conjunction of eighteen statements of the development, one for each clause; the proofs of (P1), (P6), (P8) and (H2) are given above, those of
the others are deferred (together about $12\,000$ lines): (P1) \cref{pf:sh-hodge:P1}; (P2)
\lref{GlobalTheory.NL.WachLatticeTrivialModUniformizer.wachLattice_galois_trivial_mod_uniformizer}; (P3)
\lref{GlobalTheory.NL.CmpCris.wachLattice_comparison_dcris}; (P4) \lref{GlobalTheory.T.wach_lattice_expansion_at_low_level_points};
(P5) \lref{GlobalTheory.T.wachLattice_heights_from_hodgeJumps}; (P6) \cref{pf:sh-hodge:P6}; (P7)
\lref{GlobalTheory.T.herrCocycle_poles_frobeniusQ_crystalline};
(P8) \cref{pf:sh-hodge:P8}; (P9) \lref{GlobalTheory.T.wachModule_descent_to_tame_subextension} (about $1\,950$ lines: the
construction of $\pi'$ by an Eisenstein generator of $A^\Delta$ and of a $\Delta$-invariant basis); (P10)
\lref{GlobalTheory.T.wachModule_uniqueness_characterisation}; ($\Gamma$ss) \lref{GlobalTheory.T.gammaSemisimpleQuotientD_statement}
(about $1\,300$ lines, from the eigenvalues of $\gamma$ on $Q$-power quotients of Wach modules); (H1) \lref{GlobalTheory.NL.HerrComplexH1.herr_H1_description};
(H2) \cref{pf:sh-hodge:H2}; (H3) \lref{GlobalTheory.T.wachIntegralComplex_cocycle_class_finite};
(H4) \lref{GlobalTheory.T.defect_torsion_second_cohomology_kernel} (about $2\,800$ lines); (L3)
\lref{GlobalTheory.T.cokernel_crystalline_surjection_killed_by_X_pow}; (E) \lref{GlobalTheory.NL.NygaardExtension.main} and
\lref{GlobalTheory.T.local_comparison_of_cmpOf}.
\end{proof}

\begin{corollary}[The used forms]\label{pf:sh-hodge:used}
For every datum $D$ and all Hodge--Tate, Fontaine and comparison data: at every good triple the ten clauses (P1)--(P10) hold (the
\emph{$\Delta$-invariant lattice criterion}), and ($\Gamma$ss), (H1)--(H4), (L3) and (E) hold.
\formallabel{GlobalTheory.T.wachModuleLocalConclusions_usedForms}
\formallabel{GlobalTheory.NL.usedLocalConclusion_and_prvE}
\end{corollary}

\begin{proof}
Regroup the clauses of \cref{pf:sh-hodge:wach} place by place; the four Herr clauses form the \emph{second proposition} of the development
(\lref{Formalism.MainNotions.DeltaInvariantLattices.herr_prop_two_iff}, a restatement).
\end{proof}

% ---------------------------------------------------------------------------
\subsection{Local torsion bounds}\label{pf:sh-hodge:loc}

Throughout, $D$ is a datum with Hodge--Tate, Fontaine and comparison data $(HW,F,\mathrm{Cm})$ of its local theory, $\ell$ a prime,
$(K,s,\mathit{LL})$ a good triple with Wach structure $W$ and \'etale tie $j$, $\OO=\OO_{\mathit{LL}}$, $G=G_K$, and lattices are lattices of
representations of $G$ over $\mathit{LL}$. For $a\in\Z$, $v_\ell(a)$ is the $\ell$-adic valuation ($v_\ell(0)=0$). The \emph{jump set} $J(T)$ of a
lattice $T$ is the set of $a\in\Z$ with $\mathrm{Fil}^{a+1}D_{\mathrm{dR}}(V)\ne\mathrm{Fil}^aD_{\mathrm{dR}}(V)$, the union over the embeddings
$\tau\colon K\to\mathit{LL}$ of the sets of Hodge--Tate jumps at $\tau$.

\begin{definition}[The Wach hypothesis and the exponents]\label{pf:sh-hodge:def-hypw}
Let $M_0,h_0\in\N$. A lattice $T$ satisfies $\mathrm{HypW}(M_0,h_0)$ if it is crystalline on $G_{K_s}$ (i.e.\ $\dim_{\mathit{LL}}(V\otimes
B)^{G_{K_s}}=[K_0:\Ql]\dim V$), $J(T)\subseteq[-M_0,h_0]$, and there is an $A$-submodule $N\subseteq D(T)$ with an $A$-basis whose image under
$j_T$ is an $R_W$-basis of $R_W\otimes T$, with $Q^{M_0}\varphi(N)\subseteq N$, $Q^{h_0}N\subseteq\mathrm{span}\,\varphi(N)$, stable under
$G_{K_s}$ and under all endomorphisms of $T$, and with $(g-1)N\subseteq\pi N$ for $g\in G_{K_s}$. Put
\[
n_0=\min\{n:\ M_0+h_0\le\ell^n(\ell-1)\},\qquad m=2s+n_0+1,
\]
and for a finite set $U\subseteq\Z$ and integers $a<b$
\[
c(U;a,b)=v_\ell([K_m:\Ql])-v_\ell([K_s:\Ql])+\sum_{i=a}^{b-1}\Bigl(2+\max\Bigl(1,\sum_{u\in U,\,u\ne i}\bigl(m+v_\ell(u-i)\bigr)\Bigr)\Bigr).
\]
\formal{Formalism.MainNotions.LocalSituation.HypW}
\formal{Formalism.MainNotions.LocalSituation.cLL}
\end{definition}

\subsubsection{Torsion of invariants}

\begin{theorem}[Torsion exponent of invariants]\label{pf:sh-hodge:tek}
Assume \shhodgeBK, \shhodgeLTD{} and \shhodgeLEC. Let $T$ satisfy $\mathrm{HypW}(M_0,h_0)$ and suppose that $0$ is not a Hodge--Tate jump of
$T$ at any embedding $K\to\mathit{LL}$. Then every $G_K$-invariant element $x$ of $V/T$ satisfies
\[
\ell^{\,c_{\mathrm{inv}}}x=0,\qquad c_{\mathrm{inv}}=v_\ell([K_m:\Ql])-v_\ell([K:\Ql])+2+\max\Bigl(1,\sum_{u\in J(T),\,u\ne0}\bigl(m+v_\ell(u)\bigr)\Bigr).
\]
\formallabel{GlobalTheory.T.torsion_exponent_kills_invariants_of_good}
\formallabel{GlobalTheory.HubX.InvariantsNoJump.torsion_exponent_kills_invariants_no_jump_zero}
\formallabel{GlobalTheory.HubX.InvariantsNoJump.two_add_sum_jumps_le_intervalExponent}
\formallabel{GlobalTheory.T.invariants_vanish_at_prime_of_admissible}
\formallabel{GlobalTheory.Cyc.InvKilled.invariants_killed}
\formallabel{GlobalTheory.Cyc.InvKilled.killed_quot}
\formallabel{GlobalTheory.Cyc.InvKilled.killed_of_coord}
\end{theorem}

\begin{remark}\label{pf:sh-hodge:tekrem}
The three literature hypotheses and the comparison data are not used in the proof; the hypothesis on the jump $0$ is used at one
embedding only. The result is a quantitative form of Tate's theorem $H^0(G_K,\C_\ell(j))=0$ for $j\ne0$ \cite[\S3]{Tate1967pdiv}, via
Sen's theory \cite{Sen1980}.
\end{remark}

The proof rests on two lemmas, whose proofs are deferred. Let $\C_\ell$ carry the action of $G_{\Ql}$ of the $A_{\inf}$-data and let
$\chi$ be the cyclotomic character. Choose a $\Zl$-basis of $T$ (of rank $d$) and let $M(\sigma)\in M_d(\C_\ell)$ be the matrix of
$\sigma\in G_K$ in it (\lref{GlobalTheory.T.exists_continuous_galoisLatticeMatrix_integral}). Let $K_\infty=K_0(\mu_{\ell^\infty})$,
$H_\infty=\Gal(\Omega_\ell/K_\infty)$ and $\rho_1=\ell^{-1/(\ell-1)}$. A \emph{Sen frame} of level $m$ is a map $\sigma\mapsto U(\sigma)$ from
$G_{K_m}$ to $r\times r$ matrices with integral entries in $K_m$, continuous and multiplicative, trivial on $H_\infty$, together with
$g\in G_{K_m}$ with $|\chi(g)-1|=\ell^{-m}$ and $|U(g)-1|<\rho_1$ entrywise. A \emph{Hodge--Tate basis} with weights in $J\subseteq\Z$ is a basis
$e_1,\dots,e_r$ of $\C_\ell^r$ and $u_k\in J$ with $\chi(\sigma)^{u_k}U(\sigma)\sigma(e_k)=e_k$ for $\sigma\in G_{K_m}$. Write $B_r$ for the unit ball of
$\C_\ell^r$.

\begin{lemma}[Frames]\label{pf:sh-hodge:frame}
For $T$ with $\mathrm{HypW}(M_0,h_0)$, $\tau\colon K\to\mathit{LL}$ and $m$ as above, there are $r$, a matrix $P\in M_{r\times d}(\C_\ell)$, a Sen frame
$U$ of level $m$, a finite set $J$ of Hodge--Tate jumps of $T$ at $\tau$ and a Hodge--Tate basis for $U$ with weights in $J$, such that
$PM(\sigma)=U(\sigma)\sigma(P)$ for $\sigma\in G_{K_m}$, $B_r\subseteq P(B_d)$, $\ell P(B_d)\subseteq B_r$, and every $y\in\Ql^d$ with $Py\in P(B_d)$
lies in $\Zl^d$.
\end{lemma}

This is \lref{GlobalTheory.T.exists_tau_component_frame_of_lattice} (about $1\,200$ lines; its proof uses Fontaine's bounds
\lref{GlobalTheory.Per.FontaineAll.fontaine_bounds} and the Wach module in $\mathrm{HypW}$).

\begin{lemma}[Twisted invariants]\label{pf:sh-hodge:twinv}
Let $U$ be a Sen frame of level $m\ge1$ ($m\ge2$ if $\ell=2$) with a Hodge--Tate basis with weights in a finite set $J$, let $i\notin J$, and
let $\mathcal M\subseteq\C_\ell^r$ be an additive subgroup with $B_r\subseteq\mathcal M$ and $\ell\mathcal M\subseteq B_r$. If $y\in\C_\ell^r$ satisfies
$\chi(\sigma)^iU(\sigma)\sigma(y)-y\in\mathcal M$ for all $\sigma\in G_{K_m}$, then
\[
\ell^{\,2+\max(1,\sum_{u\in J}(m+v_\ell(u-i)))}\,y\in\mathcal M .
\]
\end{lemma}

This is \lref{GlobalTheory.T.prime_pow_smul_mem_of_twistedInvariant_mod_lattice} (about $300$ lines), an integral form of the
vanishing of $H^0(G_{K_m},\C_\ell^r(i))$ for a representation with Hodge--Tate weights in $J\not\ni i$.

\begin{proof}[Proof of \cref{pf:sh-hodge:tek}]
Fix an embedding $\tau\colon K\to\mathit{LL}$ (there are $[K:\Ql]\ge1$ of them, \lref{Formalism.MainNotions.N.JsetEmb.card_emb_good}) and put
\[
c_0=2+\sum_{u\text{ jump at }\tau}\bigl(m+v_\ell(u)\bigr).
\]
\emph{Step 1: reduction to coordinates.} It suffices to show: every $y\in\Ql^d$ with $M(\sigma)y-y\in B_d$ for all $\sigma\in G_K$ satisfies
$\ell^{c_0}y\in\Zl^d$. Indeed, if $x\in V/T$ is invariant, lift it to $v\in V$; then $gv-v\in T$ for all $g$, and $v\in\ell^{-k}T$ for some $k$, so
the coordinates $y$ of $v$ satisfy the hypothesis, hence $\ell^{c_0}v\in T$, i.e.\
$\ell^{c_0}x=0$.

\emph{Step 2.} Let $r,P,U,g,J$ be as in \cref{pf:sh-hodge:frame}. We have $m\ge s+1$, so $K\subseteq K_s\subseteq K_m$ and $G_{K_m}\subseteq G_K$; and
$0\notin J$ since $0$ is not a jump at $\tau$. If $J=\emptyset$ then $r=0$ (the basis has weights in $J$), so $Py\in P(B_d)$ for every $y$ and every
$y$ in Step~1 is integral. Otherwise put $\mathcal M=P(B_d)$ and $y'=Py$. For $\sigma\in G_{K_m}$, as $y$ has rational coordinates,
\[
U(\sigma)\sigma(y')-y'=U(\sigma)\sigma(P)y-Py=P\bigl(M(\sigma)y-y\bigr)\in P(B_d)=\mathcal M .
\]
\Cref{pf:sh-hodge:twinv} with $i=0$ gives $\ell^ey'\in\mathcal M$ with $e=2+\max(1,\sum_{u\in J}(m+v_\ell(u)))$. Since $J\ne\emptyset$ the sum is at
least $1$, and it is at most the sum over all jumps at $\tau$; hence $e\le c_0$, and $\ell^{c_0}Py\in\mathcal M$ because $\mathcal M=P(B_d)$ is stable
under multiplication by $\ell$. By the last property of $P$, $\ell^{c_0}y\in\Zl^d$.

\emph{Step 3: the constant.} The jumps at $\tau$ are non-zero elements of $J(T)$, which is finite since $J(T)\subseteq[-M_0,h_0]$; so $c_0\le
2+\max(1,\sum_{u\in J(T),u\ne0}(m+v_\ell(u)))\le c_{\mathrm{inv}}$, the last because $v_\ell([K_m:\Ql])\ge v_\ell([K:\Ql])$. Hence $\ell^{c_{\mathrm{inv}}}x=
\ell^{c_{\mathrm{inv}}-c_0}\ell^{c_0}x=0$.
\end{proof}

\subsubsection{The twisted torsion class theorem}

\begin{definition}[The admitted local statements]\label{pf:sh-hodge:def-adm}
The \emph{admitted local statements} for $(D,HW,F)$ are, at every prime $\ell$ and every finite $K/\Ql$: (adm1) every $(u,w)\in B\times
B_{\mathrm{dR}}$ is $(x-\varphi x,\,x-y)$ with $x\in B$, $y\in\mathrm{Fil}^0B_{\mathrm{dR}}$; (adm2) \shhodgeLTD; (adm3) $H^2(G_K,\mathit{LL}/\OO(1))\cong
\mathit{LL}/\OO$ for every good coefficient ring; (adm4) \shhodgeLEC; (adm5) \emph{Tate's filtered invariants}: if $V$ is de~Rham and
$z\in\mathrm{Fil}^a(V\otimes B_{\mathrm{dR}})$ satisfies $gz-z\in\mathrm{Fil}^b$ for all $g$, then $z\in(\mathrm{Fil}^aD_{\mathrm{dR}}(V))+\mathrm{Fil}^b$;
(adm6) continuity of cocycles defined by periods; (adm7) Fontaine's bounds $\dim D_{\mathrm{cris}}(V)\le[K_0:\Ql]\dim V$, $\dim D_{\mathrm{dR}}(V)\le
[K:\Ql]\dim V$, and crystalline $\Rightarrow$ de~Rham; (adm8) duals of de~Rham representations are de~Rham, and $\dim\mathrm{Fil}^0D_{\mathrm{dR}}(V)
+\dim\mathrm{Fil}^0D_{\mathrm{dR}}(V')=\dim D_{\mathrm{dR}}(V)$ for twisted duals; (adm9) at every good triple, Fontaine's statements on the \'etale tie
($\varphi-1$ is onto $R_W\otimes T$, its kernel is $1\otimes T$, cocycles of the cyclotomic kernel come from $R_W\otimes T$) and the
decompletion statement for $j$.
\formal{Formalism.MainNotions.LocalSituation.admitted}
\end{definition}

The development proves the admitted local statements, and ties, from \shhodgeLTD{} and \shhodgeLEC{} for every datum and all $HW,F$
(\lref{GlobalTheory.LocInputs.ties_and_inputs_of_duality_euler}); only (adm2) and (adm4) use these hypotheses.

\begin{theorem}[Theorem TB]\label{pf:sh-hodge:tb}
Assume \shhodgeBK, \shhodgeLTD{} and \shhodgeLEC. For every datum $D$, all data $(HW,F,\mathrm{Cm})$, every prime $\ell$ and every good triple
$(K,s,\mathit{LL})$ the following hold, where $e\in\N$ and $V_e$ denotes a representation of dimension one on which $g\in G$ acts by
$\chi(g)^{-e}$, with a lattice $L_e$.
\begin{enumerate}[label=(\alph*)]
\item \emph{(Extension classes.)} Let $0\to Y\xrightarrow{\iota}\tilde Y\to Z\to0$ be a short exact sequence of lattices with $Z$ of rank one
  and trivial action, $\tilde Y$ crystalline on $G_{K_s}$ and $Y$ satisfying $\mathrm{HypW}(M_0,h_0)$; let $\tilde y\in\tilde Y$ map to a generator
  of $Z$, $f\colon G\to Y$ with $\iota(f(g))=g\tilde y-\tilde y$, and $\phi\colon Y\to V_e/L_e$ an $\OO$-linear $G$-map. Then
  \[
  \ell^{\,c_{TB}}\,(\phi\circ f)\in\shhodgeHfD(L_e),
  \]
  \[
  c_{TB}=c(J(Y);0,e+1)+(e+2)s+v_\ell(e)+\lfloor e/(\ell-1)\rfloor+1+s+v_\ell(e+1)+v_\ell([K_s:\Ql])-v_\ell([K:\Ql]).
  \]
\item \emph{(Rank one twists.)} For every $n$, the $\ell^n$-torsion of $H^1(G,V_e/L_e)/\shhodgeHfD(L_e)$ (locally constant cocycles modulo the
  finite ones) has $\OO$-length at most $([K:\Ql]+1)\,n\,\len_\OO(\OO/\ell)$; if $e>0$ then $\shhodgeHfD(L_e)=B^1(V_e/L_e)$; and if $e=0$, a cochain
  $c$ lies in $\shhodgeHfD(L_0)$ iff it is a locally constant cocycle whose restriction to the inertia subgroup of $G_K$ is a coboundary.
\end{enumerate}
\formallabel{GlobalTheory.T.theorem_tb_of_local_blochKato_duality_euler}
\formallabel{GlobalTheory.TheoremTBRecords.aux_of_blochKato_tateDuality_eulerChar}
\formallabel{GlobalTheory.T.twistedTorsionClassTheorem_goodCase}
\formallabel{GlobalTheory.TheoremTB.tb_assembled}
\formallabel{GlobalTheory.T.weight_good_reduction_at_prime}
\formallabel{GlobalTheory.TwistTors.twist_torsion_facts}
\formallabel{GlobalTheory.TwistTors.tbe_of}
\formallabel{GlobalTheory.TwistTors.exists_unr}
\formallabel{GlobalTheory.T.crystalline_extension_class_bound}
\formallabel{GlobalTheory.Cyc.ClassBound.class_bound}
\end{theorem}

\begin{remark}
In the development (a) is the statement $\mathrm{TB}$ and (b) the statement $\mathrm{TBe}$ at the good triple (together \emph{Conclusion~4}).
Part (b) for $e=0$ is the familiar identification of $H^1_f(K,\mathit{LL}/\OO)$ with the unramified classes, and for $e>0$ the vanishing of
$H^1_f(K,\mathit{LL}(-e)/\OO(-e))$ \cite[\S3]{BlochKato1990}; part (a) bounds the failure of an integral extension class to be finite in terms
of the Hodge--Tate jumps.
\end{remark}

\begin{proof}
By \lref{GlobalTheory.LocInputs.ties_and_inputs_of_duality_euler} the admitted local statements hold for $(D,HW,F)$, and by
\cref{pf:sh-hodge:bkall} the Bloch--Kato statements $\mathrm{bk}(\mathrm{PD}^D_K)$ hold for all $\ell,K$.

\emph{Part (b).} $V_e$ is de~Rham with $D_{\mathrm{dR}}(V_e)\subseteq\mathrm{Fil}^0$
(\lref{GlobalTheory.isDR_and_DdR_le_Fil_zero_of_twist}). If $e>0$ there is $g\in G_K$ with $\chi(g)^e\ne1$
(\lref{GlobalTheory.Cyc.exists_mem_GK_chi_pow_ne_one}), so $V_e^{G}=0$, and by (adm5) and (adm7) every finite cocycle of $V_e$ is a coboundary
(\lref{GlobalTheory.HfZero.z1f_le_b1}); hence $\shhodgeHfD(L_e)=B^1(V_e/L_e)$, which is non-zero. If $e=0$, $V_0$ is the trivial representation and (Ad)
gives $\dim\shhodgeHf(V_0)=1+(1-1)=1$. Choose an unramified continuous homomorphism $\delta\colon G_K\to\Ql$ which is not zero; it is a finite cocycle of the
trivial representation (\lref{GlobalTheory.UnramZl.exists_unramified_character} and \lref{GlobalTheory.UnramifiedCrys.unramified_mem_Z1f}). Since $G_K$ is
topologically generated by a Frobenius lift and the inertia group (\lref{GlobalTheory.exists_topologicalClosure_zpowers_sup_inertia}),
\lref{GlobalTheory.UnramifiedClasses.HfD_eq_unramified} identifies $\shhodgeHfD(L_0)$ with the locally constant cocycles that are coboundaries on
inertia. Finally (adm2) and (adm4) bound the number of locally constant cocycles of a finite module $M$ killed by $\ell^r$ by $(\#M)^{[K:\Ql]+2}$
(\lref{GlobalTheory.CocycleCount.cocycles_finite_ncard_le}); with $\shhodgeHfD(L_e)\ne0$ this gives the length bound for the $\ell^n$-torsion
(\lref{GlobalTheory.TwistTors.len_part}), using $\len_\OO(\OO/\ell^n)=n\len_\OO(\OO/\ell)$.

\emph{Part (a)} is the \emph{crystalline extension class bound}, which uses only the admitted statements. Its proof proceeds as follows (the three main
steps are deferred). Let $\xi\in A_{\inf}$ be the standard generator of $\ker\theta$ and $Q_{a,b}=\xi^aA_{\inf}/\xi^bA_{\inf}$.
\begin{enumerate}[label=(\arabic*)]
\item $f$ is a continuous cocycle, and $f\otimes1$ is a cocycle of $G_{K_s}$ in $Y\otimes_{\Zl}Q_{0,e+1}$.
\item Since $\tilde Y$ is crystalline on $G_{K_s}$, $\ell^k(f\otimes1)$ is a coboundary for some $k$
  (\lref{GlobalTheory.T.ses_lattice_trivial_quotient}).
\item The $\ell$-power torsion of the first cohomology of $G_{K_s}$ in $Y\otimes Q_{a,b}$ (computed with continuous cocycles) is killed by
  $\ell^{c(J(Y);a,b)}$ when $Y$ satisfies $\mathrm{HypW}(M_0,h_0)$ (\lref{GlobalTheory.T.torsion_bound_of_weight_hyp}, about $2\,500$ lines);
  hence $\ell^{c(J(Y);0,e+1)}(f\otimes1)$ is a coboundary.
\item If $\ell^N(f\otimes1)$ is a coboundary in $Y\otimes Q_{0,e+1}$, then $\ell^{N+c'}\phi\circ f\in\shhodgeHfD(L_e)$ with $c'=c_{TB}-c(J(Y);0,e+1)$, by
  (adm2) and (adm4) (\lref{GlobalTheory.Cyc.ClassBound.mem_HfD_of_lattice}, about $5\,700$ lines).
\end{enumerate}
\end{proof}

\subsubsection{The local lifting theory}

Let $W$ be a Wach structure, $T_1,T_2$ crystalline lattices and $F\colon T_1\to T_2$ a surjective lattice map, and let $T_1$ have jumps within
$(M,h)$. Put $M_{\max}=\max_iM_i$, $h_{\max}=\max_ih_i$, $H=\max_i(M_i+h_i)$, let $J_1$ be the set of jumps of $T_1$ in $[-M_{\max},h_{\max}]$ and
\[
J(J_1,M)=\{a+j\ne0:\ a\in J_1,\ 1\le j\le M_{\max}\}\cup[1,M_{\max}],\qquad c_H=\lceil1+\log_\ell(H+1)\rceil,
\]
\[
e_{WA}=\Bigl\lfloor\frac{H\ell}{(\ell-1)^2}+3+\log_\ell(H+1)\Bigr\rfloor,\qquad a_A=\Bigl\lfloor\frac{M_{\max}}{\ell-1}\Bigr\rfloor+1,
\]
\[
e_{PA}=2e_{WA}+\sum_{j\in J(J_1,M)}v_\ell(j)+a_A+v_\ell(a_A!).
\]
$s_i$ denotes the least $k$ with $\pi^kN_i(T_2)\subseteq F(N_i(T_1))$.
\formal{Formalism.MainNotions.LiftingTheory.ePA}

\begin{theorem}[The local constants]\label{pf:sh-hodge:lc}
Assume \shhodgeBK, \shhodgeLTD{} and \shhodgeLEC. For every datum $D$ and all data $(HW,F,\mathrm{Cm})$, the \emph{local lifting theory with
constants $(C_R,C_A,C_B,C_K)=(10,23,32,32)$} holds at every good triple: that is, the following eight packages.
\begin{enumerate}[label=(\arabic*)]
\item (L) An ideal statement: for disjoint finite $S_1,S_2\subseteq\Z$ and $c\colon\Z\to\OO$ with $v(c(a)-c(b))=s_g+v_\ell(a-b)$ for $a\ne b$ in
  $S_1\cup S_2$, if $v(\prod_{b\in S_2}(c(a)-c(b)))+v(f_1'(c(a)))\le E$ for all $a\in S_1$, where $f_1=\prod_{b\in S_1}(X-c(b))$, then $p^E$ lies in
  the ideal of $\OO[X]$ generated by $f_1$ and $\prod_{b\in S_2}(X-c(b))$.
\item (WX) Wach modules kill cokernels: $X^H$ kills $N_i(T_2)/F(N_i(T_1))$ for every letter $i$; in case~A so does $p^{e_{WA}}$; in case~B
  ($G_{K_s}=G_K$, $[K:K_0]=(\ell-1)\ell^{s-1}$) so does $p^{2Hs+v_\ell((2H)!)}$; and if $s=1$, $\ell\ne2$, $[K:K_0]=e$, a $\Delta$-statement with
  $\pi'$ and the corresponding power of $p$.
\item (W1) In case~A: $p^a\pi^j$ kills $N_i(T_2)/F(N_i(T_1))$ whenever $(s_i-j)\ell/(\ell-1)^2+c_H\le a$.
\item (WR) In case~A: $s_{i-1}\le M_i+h_i+c_H+3+s_i/(\ell-1)$.
\item (WC) In case~A, with $C_R=10$: explicit bounds for $s_i$, for the least $r_i$ with $p^{r_i}$ killing the $\Delta$-invariants, and for
  $\lceil s_i/(\ell-1)\rceil+r_i$, among them $r_i\le C_R\bigl(G/\ell+H/(\ell-1)^{j+1}+1+\log(H+1)\bigr)$ when $M_{i+a}+h_{i+a}\le2G$ for
  $1\le a\le j$.
\item (H) In case~A: the action of $\gamma$ on $\pi_0^r$-multiples of $\Delta$-invariants, the valuations $v(u_1^r-1)=1+v_\ell(r)$ of
  $u_1=\chi(\gamma)^{\ell-1}$, and bounds for the solutions of the associated recursions in a $\Delta$-frame.
\item (Tr) In case~A: decompositions of $\varphi^k(\pi_0+p)$ and of $\varphi^j(\pi_0)^r$ in $A^\Delta$, with bounds for the denominators.
\item (Lift) The lifting theorems: in case~A, $F$ lifts with exponent $e_{PA}$ (i.e.\ $\ell^{e_{PA}}\shhodgeZf(T_2)$ lies in the image of $\shhodgeZf(T_1)$
  plus coboundaries, \cref{pf:linkN-adm:def-col}); in case~B with the analogous exponent $e_{PB}$; for crystalline $T_1$ with jumps within $(M,h)$,
  the product $\prod_{j\in J(J_1,M)}(1-u^{-j})$ is associated to a power of $p$ in case~B, and $\prod_{j\in J(J_1,M)}\frac{u^{-j}-1}{p}=wp^k$ with
  $w\in\OO^\times$, $k\in\Z$ in case~A, where $u=\chi(\gamma)$ as in \cref{pf:sh-hodge:hw}; the exponent bounds
  \[
  e_{PA}\le C_A\Bigl(1+\frac{H}{\ell}+1+\log_\ell(H+1)\Bigr)\ \text{in case~A},\qquad e_{PB}\le C_B(Hs+1),\quad e_{PB}\le C_KH\ \text{if }s=1,
  \]
  in case~B; and Galois descent of the lifting theorems to subfields $K'\subseteq K$ with $G_K$ normal in $G_{K'}$.
\end{enumerate}
\formallabel{GlobalTheory.T.localConstants_from_admitted_local_Galois_statements}
\formallabel{GlobalTheory.localConstants_from_admitted_local_Galois_statements_alt}
\formallabel{Formalism.MainNotions.LiftingTheory.frobenius_lifting}
\formallabel{Formalism.MainNotions.LiftingTheory.local_theory_of_admitted}
\end{theorem}

\begin{proof}
The development defines twenty-three \emph{admitted hypotheses of the lifting theory}: the clauses (P1)--(P10), ($\Gamma$ss), (H1)--(H4) and (L3)
of \cref{pf:sh-hodge:hw}, and seven further local statements (the Herr cocycle of an extension of the unit object, also in its $\Delta$-form; a
descending sequence for $\gamma$ at the jumps; corestriction for subfields; the naturality of $\shhodgecl$ and $\shhodgecl_K$; Fontaine's
inequality; and the exactness of Hodge--Tate jumps in short exact sequences). They hold for all data
(\lref{GlobalTheory.T.liftingTheorem_wachModuleLemma_hypotheses_admitted}). From them, \lref{Formalism.MainNotions.LiftingTheory.local_parts_of_admitted}
(deferred, about $1\,600$ lines) proves the packages (1)--(7), and \lref{Formalism.MainNotions.LiftingTheory.lifting_all_parts_of_admitted}
(deferred, about $4\,500$ lines) proves (8): the case~A lifting theorem \lref{Formalism.MainNotions.LiftingTheory.lifting_case_A_of_admitted}, the
case~B theorem \lref{Formalism.MainNotions.LiftingTheory.lifting_case_B_of_hypotheses}, the descent to subfields
\lref{Formalism.MainNotions.LiftingTheory.lifting_descent_to_subfields} and the exponent shapes
\lref{Formalism.MainNotions.LiftingTheory.lifting_exponent_shapes}.
\end{proof}

\begin{remark}
The hypotheses \shhodgeBK, \shhodgeLTD{} and \shhodgeLEC{} of \cref{pf:sh-hodge:lc} are not used in its proof: the admitted hypotheses of the
lifting theory are proved for all data without them.
\end{remark}

% ---------------------------------------------------------------------------
\subsection{Three global statements}\label{pf:sh-hodge:glob}

\subsubsection{The parameters of a family grow polynomially}

We use the notation of \cref{sec:pf:linkN,sec:pf:linkN-sel} for families: a family $\mathcal F\in\sfname{FamP}(\eps)$ has index sets $I_x$, a
threshold $x_1$, members $D_{x,i}$ with size data $L_{x,i}$ of scale $(x,\eps)$ ($x>1$, $0<\eps<1/52$, $\Lambda=x^\eps$, $b=e^\Lambda$, and a
size parameter $\delta_2>0$), a constant $C_t$, a list of constants from which the parameter choices of the members are read off (in particular
two constants $C_{Lw},C_{YL}$), and the \emph{size conditions} of \cref{sec:families:galois}. Among them we use:
\begin{equation}\label{eq:pf:sh-hodge:sizes}
\begin{gathered}
n\le x^{2+8\eps+\eps/100},\quad k\le x^{9/2+10\eps+\eps/100},\quad 1\le t\le x^{C_t+\eps/100},\quad x^{-\eps}\le\delta_2,\\
\log q,\ \log d\le x^{\eps/5+\eps/100},\quad 1\le\mathrm{rd}_E\le\tfrac{Q_S}{2}\sqrt{qd},\quad (kb)^8\le g_j,\quad \log g_j\le x^{\eps/100}\Lambda\quad(j=0,1),
\end{gathered}
\end{equation}
where $\mathrm{rd}_E=|d_E|^{1/[E:\Q]}$ is the root discriminant of $E$ and $Q_S\le n^2$ is the level of the host. Recall $f=2n$, $r=k+m-1$ and
$m+3\le k$, so $r\le2k$; $\#G$ is the order \eqref{eq:pf:linkN:G} of the level group.

\begin{theorem}[Parameter bound]\label{pf:sh-hodge:par}
For every $\eps$ and every family $\mathcal F\in\sfname{FamP}(\eps)$ there are $C\ge1$ and $x_0$ such that for all $x\ge x_0$ with $x>x_1$ and all
$i\in I_x$ the member $(D,L)=(D_{x,i},L_{x,i})$ satisfies
\[
C_{Lw}\le C,\quad C_{YL}\le C,\quad k,n,t\le x^C,\quad x^{-\eps}\le\delta_2,\quad \log h_E\le Cx^C,\quad \log\#G\le Cx^C\Lambda .
\]
\formallabel{GlobalTheory.T.admissible_parameters_family_const}
\formallabel{GlobalTheory.NLink.admissible_parameters_design_family}
\end{theorem}

\begin{lemma}[Class number]\label{pf:sh-hodge:hE}
For every $\eps$ and $\mathcal F\in\sfname{FamP}(\eps)$, with $A=|2+8\eps+\eps/100|+|\eps/5+\eps/100|+1$, every member with $x\ge1$, $x>x_1$ satisfies
$\log h_E\le(30+3A)x^{30+3A}$.
\formallabel{GlobalTheory.T.log_hostClassNumber_le_scale_pow_family}
\formallabel{GlobalTheory.log_hE_le_pow}
\formallabel{GlobalTheory.ClassNumHost.log_hE_le_lin}
\formallabel{GlobalTheory.ClassNumHost.log_rd_le}
\formallabel{GlobalTheory.ClassNumHost.log_hE_le_rd}
\end{lemma}

\begin{proof}
$E$ is totally complex of degree $2f$, and \lref{GlobalTheory.log_classNumber_le_log_abs_discr_div} states that a totally complex field of
degree $2f$ satisfies $\log h\le f\log\mathrm{rd}+2f\log(1+f\log\mathrm{rd})$ with $\log\mathrm{rd}=\log|d|/(2f)$. Since $\mathrm{rd}_E\ge1$ and
$\log(1+y)\le y$ for $y\ge0$,
\[
\log h_E\le(1+2f)\,f\log\mathrm{rd}_E .
\]
From $\mathrm{rd}_E\le\frac{Q_S}2\sqrt{qd}\le n^2\sqrt{qd}$ and $\log n\le n$ we get $\log\mathrm{rd}_E\le2n+\frac12(\log q+\log d)$. Put $P=x^A\ge1$.
By \eqref{eq:pf:sh-hodge:sizes}, $n\le P$ and $\log q,\log d\le P$ (the exponents are at most $A$ and $x\ge1$); hence $\log\mathrm{rd}_E\le3P$, $f=2n\le2P$,
$f\log\mathrm{rd}_E\le6P^2$ and $1+2f\le5P$. Therefore $\log h_E\le30P^3=30x^{3A}\le(30+3A)x^{30+3A}$.
\end{proof}

\begin{lemma}[Level group]\label{pf:sh-hodge:G}
For every $\eps$ and $\mathcal F\in\sfname{FamP}(\eps)$ there is $x_0$ such that every member with $x\ge x_0$, $x>x_1$ satisfies $\log\#G\le x^{15}$,
$8fr/\min(g_0,g_1)\le1/3$ and $b\ge1$; in particular $\log\#G\le15x^{15}\Lambda$.
\formallabel{GlobalTheory.T.exists_log_card_restRows_le_pow_mul_length}
\formallabel{GlobalTheory.NLink.log_card_levelGroup_le_scale_pow_mul_length}
\formallabel{GlobalTheory.log_levelGroup_le_and_level_large}
\end{lemma}

\begin{proof}
By the size conditions there is $x_0'$ such that for $x\ge x_0'$, $x>x_1$, every member satisfies $f\le x^3$, $k\le x^5$ and $x^{C_t+6}\le b$, hence $b\ge1$
(\lref{GlobalTheory.poly_sizes_ranks_weights}; for instance $f=2n\le2x^{2+8\eps+\eps/100}\le x^{1/2}x^{5/2}$ when $x\ge4$, as $\eps<1/52$).
Let $x_0=\max(x_0',8)$. By \eqref{eq:pf:sh-hodge:sizes}, $g_j\ge(kb)^8\ge k^8>0$ and $\log g_j\le x^{\eps/100+\eps}\le x$. Since
$\#\GL_r(\F_{g})\le g^{r^2}$, \eqref{eq:pf:linkN:G} gives $\#G\le(g_0g_1)^{r^2f}$, so
\[
\log\#G\le r^2f(\log g_0+\log g_1)\le(2x^5)^2\,x^3\,(2x)=8x^{14}\le x^{15}.
\]
Next $8fr\le8k\cdot2k=16k^2$ (as $f\le k$) and $\min(g_0,g_1)\ge k^8\ge5^6k^2$ (as $k\ge m+3\ge5$), so $8fr/\min(g_0,g_1)\le16/5^6\le1/3$. Finally
$\log\#G\le x^{15}\le15x^{15}\le15x^{15}\Lambda$ since $\Lambda=x^\eps\ge1$.
\end{proof}

\begin{proof}[Proof of \cref{pf:sh-hodge:par}]
Let $C_1=30+3A$ and $y_1=1$ (\cref{pf:sh-hodge:hE}), $C_2=15$ and $y_2$ the threshold of \cref{pf:sh-hodge:G}, and
\[
C=\max\bigl(1,\ C_{Lw}(\mathcal F),\ C_{YL}(\mathcal F),\ 9/2+10\eps+\eps/100,\ 2+8\eps+\eps/100,\ C_t+\eps/100,\ C_1,\ C_2\bigr),
\]
where $C_{Lw}(\mathcal F),C_{YL}(\mathcal F)$ are the entries of the constant list of $\mathcal F$ which are the constants $C_{Lw},C_{YL}$ of every
member's parameter choice. Let $x\ge x_0=\max(y_1,y_2)$, $x>x_1$; then $x\ge1$, and $x^a\le x^C$ for $a\le C$. The first two bounds hold by the choice
of $C$; $k,n,t\le x^C$ and $x^{-\eps}\le\delta_2$ follow from \eqref{eq:pf:sh-hodge:sizes}; $\log h_E\le C_1x^{C_1}\le Cx^C$ and
$\log\#G\le C_2x^{C_2}\Lambda\le Cx^C\Lambda$ by the lemmas, using $Cx^C\ge C_ix^{C_i}$ for $x\ge1$, $0\le C_i\le C$.
\end{proof}

\subsubsection{The signature choice condition}

Let $D$ be a datum with a norm choice $\mathcal C$ (\cref{sec:families:galois}): for every twisted place $w$ of the coefficient field it provides a
twisting character $\psi_w$ of a ray class group with values in a finite extension of $\Z_{\ell(w)}$, a side $\mathrm{sd}(w)$, an exponent $k'(w)$
and a finite set $\mathcal Q(w)$ of auxiliary places. For a \emph{signature} $s$ of the twist at $w$, $\mathrm{sLen}(w,s)\in[0,\infty]$ denotes the
length of the Selmer group of the Bloch--Kato reduction of the twist for $s$, relaxed at $\mathcal Q(w)$, for the character $\psi_w$ read at the side of
$s$; and $\mathrm{tv}_w(P)$ the tame valuation of $\psi_w$ at $P$ (\cref{pf:linkN-sel:kill}).

\begin{definition}[Signature choice condition]\label{pf:sh-hodge:def-sig}
Let $\rho=(\rho_\iota)$ be a family of rank $r$. The \emph{signature choice condition} for $\rho$ and $\mathcal C$ holds if for every twisted local
index $\iota$, with $w=w_\iota$: (a) for every $\pi\in Z^{\mathrm{good}}_\iota$ there is a good signature $s$ of the twist at $w$ such that
$\mathrm{sLen}(w,s)<\infty$ implies that $\shhodgeSel(\iota,\rho,\pi)$ is finite, and $s_\Pi(\iota,\rho,\pi)\le\mathrm{sLen}(w,s)$; and (b)
$\sum_{v\in Q_\iota}\lambda_\iota(\OO_\iota/(\zeta_{\iota,v}-1))\le\sum_{P\in\mathcal Q(w)}\mathrm{tv}_w(P)$.
\formal{Formalism.MainNotions.N.KillingConditions.SignatureChoiceCondition}
\end{definition}

The Chebotarev statement used here, $\sfname{Cheb}$, is: for every number field $F$ with algebraic closure $\Omega$, every open normal subgroup
$U\subseteq G_F$, every $g\in G_F$ and every finite set $A$ of finite places, there are $v\notin A$ of absolute residue degree and ramification index
$1$, a place $B$ above $v$ with $\shhodgeIne_B\subseteq U$ and a Frobenius $\sigma$ at $B$ for $\shhodgeNm v$ with $\sigma g^{-1}\in U$. It is equivalent to
the statement of \cref{pf:sh-hodge:cheb} (\lref{GlobalTheory.Published.chebotarevDegreeOne_iff_cheb}), and holds by that theorem
(\lref{GlobalTheory.Published.cheb_holds}).
\formal{GlobalTheory.TwSys.Cheb}

\begin{theorem}[Signature choice]\label{pf:sh-hodge:js}
Assume $\sfname{Cheb}$. Let $D$ be a datum with norm choice $\mathcal C$ and families $r,\rho$ of ranks $n_V,r$. Assume that $(r,\rho)$ is a
\emph{point family}, that $\mathcal C$ is \emph{tied} to $D$, and the three conditions of the norm-choice facts on good triples, on the same twist and
on the good-point rows, and residual absolute irreducibility at good points. Then the signature choice condition holds for $\rho$ and $\mathcal C$.
\formallabel{GlobalTheory.JointSix.main}
\formallabel{GlobalTheory.PtSel.selCore}
\formallabel{GlobalTheory.JointSix.psiAt_psiW}
\end{theorem}

\begin{proof}
Let $\iota$ be twisted and $w=w_\iota$. Since $\mathcal C$ is tied to $D$, the twisting index of $\iota$ is the index of $w$ and $\mathcal Q(w)=Q_\iota$.
\emph{(a)} Let $\pi\in Z^{\mathrm{good}}_\iota$. The point family gives that $\rho_{\iota,\pi}$ is a point representation at $(\iota,\pi)$, and the same-twist
condition gives a character $\psi$ and a ring map identifying the twisting character of $\iota$ with $\psi$ (the \emph{same character} condition). The
statement \lref{GlobalTheory.PtSel.sel_at} (deferred, about $2\,060$ lines) provides a good signature $s$ whose side is the twist side of $\iota$,
such that $\mathrm{sLen}<\infty$ implies finiteness of $\shhodgeSel(\iota,\rho,\pi)$ and $s_\Pi(\iota,\rho,\pi)\le\mathrm{sLen}$, where $\mathrm{sLen}$ is
computed with $\psi$ and $Q_\iota$. Its proof: the good-point row data at $(\iota,\pi)$ and the good-triple condition provide $s$; by $\sfname{Cheb}$ and
residual absolute irreducibility the twisted $\rho_{\iota,\pi}$ is conjugate to the representation on the lattice attached to $\psi$
(\lref{GlobalTheory.PtCore.twistedRho_conj_latTpsi}, using that representations over a discrete valuation ring with the same traces and absolutely
irreducible residual representation are conjugate, \lref{GlobalTheory.T.exists_conj_of_trace_eq_residually_absIrreducible_DVR}); Selmer lengths are
invariant under this conjugation (\lref{GlobalTheory.ConjLatticeSelmer.sPiTop_eq_sTLoc}) and under the change of character
(\lref{GlobalTheory.SelmerSameChar.sTLoc_eq}); and the Selmer group of the twist is compared with that of the Bloch--Kato reduction
(\lref{GlobalTheory.SelTwistLe.sT_le_sLen}). Finally the character of the side of $s$ is $\psi_w$: reading $\psi_w$ at the side
$\mathrm{sd}(w)$ gives the character $\psi_G$ used above, by inspection of the two cases of the side.
\emph{(b)} For $v\in Q_\iota$, $\lambda_\iota(\OO_\iota/(\zeta_{\iota,v}-1))\le\mathrm{tv}_w(v)$ by
\lref{GlobalTheory.T.conductorRoot_sub_one_length_le_tameValue} (deferred); sum over $v\in Q_\iota=\mathcal Q(w)$.
\end{proof}

\subsubsection{The global lattice index theorem}

We use the notation of \cref{pf:linkN-sel:glit}. Let $D$ be a datum, $\iota$ a local index with $\ell=\ell_\iota$, $r$ and $\rho$ families of ranks $n_V$ and
$r$ at $\iota$, $\pi\in X_\iota$, $\mathrm{dc}=(\mathrm{row},\mathrm{col})$ a partition of the places of $E$ above $\ell$, $\kappa_{\mathrm{row}}$ a rational
function on places, $e'$ an integer function on places and $c_B\in\N$. The \emph{admitted conditions} at $(\iota,r,\rho,\pi,\mathrm{dc},\kappa_{\mathrm{row}},e',c_B)$ are the nineteen conditions of \cref{pf:linkN-sel:glit}:
(G1) the coefficients of the characteristic polynomials of the $r_x(g)$, as functions of $x\in X_\iota$, lie in $T_\iota$; (G2) at the places above $\ell$, a
gain bound and crystallinity of $r_x\chi_{\iota,1}^{-1}$; (G3) absolute irreducibility of the $r_x$, the residual characteristic polynomial of $r_x$ is that
of $\rho_\pi$ times $\prod_{c=1}^t(X-\bar\chi_{\iota,c})$, and the residual image of $\rho_\pi$ spans the matrix algebra; (G4) unramifiedness outside $\ell$,
$S_\iota$, $Q_\iota$; (G5) conditions on inertia at $S_\iota\cup Q_\iota$; (G6) if $r_x$ has an eigenvector with character $\chi_{\iota,1}$ then $x\in Z_\iota$;
the tame generator conditions; local Bloch--Kato statements at the places above $\ell$ (top-degree duality, dimension formula, integral duality) and the
Bloch--Kato annihilators there (\cref{pf:sh-hodge:bka}); continuity of the family; the semisimple decomposition $r_\pi\otimes K_\iota\cong\rho_\pi\oplus
\bigoplus_{c=1}^t\chi_{\iota,c}$; Fontaine's statements and the finiteness of local cohomology at the places above $\ell$ (\cref{pf:sh-hodge:fin}); the
compatibility of the local coefficients; a condition (G51) on $(r,\rho,\pi)$; the triviality of unramified finite classes at the columns; the row length bounds
with $\kappa_{\mathrm{row}}$; the lift exponent conditions at the columns with $e'$; and the block condition with $c_B$.
\formal{Formalism.MainNotions.TruncatedSeriesCountingIdentities.GlobalLatticeIndexTheorem}

\begin{theorem}[Global lattice index theorem]\label{pf:sh-hodge:glit}
Assume \shhodgeGAR, \shhodgeLTD{} and \shhodgeLEC, and let the data above satisfy the admitted conditions. If $\ell\notin\{g_0,g_1\}$, $\pi\in Z_\iota$,
$\shhodgeSel(\iota,\rho,\pi)$ is finite and $\Sha^1(G_{E,S_\ell},\Ql/\Zl(-e))$ is finite for $0<e<t$, then inequality \eqref{eq:pf:linkN-sel:glit}
holds:
\begin{multline*}
v_{cc}(\iota,\pi)\le\sum_{w\in\mathrm{row}\cup\mathrm{col}}\kappa(w)+s_\Pi(\iota,\rho,\pi)+\sum_{0<e<t}\mathrm{sh}(\ell,e)+v_\ell(h_E)
+t\sum_{w\in\mathrm{row}}(n_w+2)c_B\\
+4n_V(\#\mathrm{col}+1)\delta_\iota+t\,\#\mathrm{col}\cdot h_0(\iota,\mathrm{col})+t\sum_{v\in Q_\iota}\lambda_\iota(\OO_\iota/(\zeta_{\iota,v}-1)),
\end{multline*}
where $\kappa(w)=\kappa_{\mathrm{row}}(w)$ for $w\in\mathrm{row}$ and $\kappa(w)$ is the column length defined from $e'$ for $w\in\mathrm{col}$.
\formallabel{GlobalTheory.T.globalLatticeIndexTheorem_of_lengthBoundAtPlaces}
\formallabel{GlobalTheory.SelmerBoundAtPoint.main}
\end{theorem}

\begin{proof}
The proof assembles four results, whose proofs are deferred.
\begin{enumerate}[label=(\arabic*)]
\item \emph{Local duality and class field facts} (\lref{GlobalTheory.T.local_duality_class_field_facts_of_artin_tate_euler}, about $7\,200$ lines): from
  \shhodgeGAR, \shhodgeLTD{} and \shhodgeLEC{} the development derives thirteen statements used in the argument
  (\lref{GlobalTheory.langRecCore_of_artinReciprocity}): the identification of decomposition groups with local Galois groups; from \shhodgeLTD{}
  the invariant isomorphisms, also for cyclotomic twists, and the finite, rational and integral forms of local duality for dual lattices
  (\lref{GlobalTheory.T.invIso_of_localTateDuality}, \lref{GlobalTheory.T.natInvIso_of_localTateDuality_cyclotomicTwist}, and \lref{GlobalTheory.T.integral_localDuality_of_localTateDuality}); from \shhodgeLTD{} and \shhodgeLEC{} the shape of local cohomology and of
  homomorphisms modulo unramified ones (\lref{GlobalTheory.T.cohomology_shape_of_tate_duality_euler} and \lref{GlobalTheory.T.galoisHom_modulo_unramified_shape_of_localDuality}); from \shhodgeGAR{} the bound of the unramified homomorphisms of $G_E$
  by the class number, which produces the term $v_\ell(h_E)$ (\lref{GlobalTheory.T.unramifiedGaloisHoms_bounded_classNumber_of_artinReciprocity});
  and unramified decompositions and elementary facts on local cohomology and topologies
  (\lref{GlobalTheory.T.unramified_padic_decomposition_hom_exists}).
\item \emph{Finiteness of twisted invariants} (\lref{GlobalTheory.SelmerBoundAtPoint.finCore}): at the places above $\ell$ the invariants of the
  twisted coefficient modules $K_\iota/\OO_\iota$ attached to $\rho_\pi$ and to the negative Tate twists have finite length.
\item \emph{The unramified input} (\lref{Formalism.MainNotions.N.K.StD.unramified_input_of_admitted}), a reformulation of the triviality of unramified
  finite classes at the columns.
\item \emph{The theorem from the inputs} (\lref{Formalism.MainNotions.N.K.Closed.theorem_n5_from_inputs}, about $26\,500$ lines), which derives the
  inequality from the admitted conditions and (1)--(3). Its proof first replaces $r_\pi$ by the block diagonal representation $\rho_\pi\oplus\bigoplus_c
  \chi_{\iota,c}$ (permitted by the semisimple decomposition; \lref{Formalism.MainNotions.N.K.Norm.normAt_rN}); then establishes, for the normalised
  family, five \emph{open inputs} (a statement on $\rho$, the block statement with $c_B$, a statement on the characters $\chi_{\iota,c}$, the first local
  condition and the row tie-in; \lref{Formalism.MainNotions.N.K.Closed.open_inputs_at_normalized}); and finally derives the inequality from ring data
  attached to the normalised family, in which the length of $\OO_\iota/\mathfrak c_\pi$ is compared with lengths of Selmer groups of the blocks
  (\lref{Formalism.MainNotions.N.K.normalized_ring_data}, \lref{Formalism.MainNotions.N.K.St.selmer_inequality_three}, and \lref{Formalism.MainNotions.N.K.Final.theorem_n5_of_open_b}).
\end{enumerate}
Given (1) and (2), the theorem is (4) applied to the admitted conditions, with the unramified input (3) in place of the column triviality.
\end{proof}

% ---------------------------------------------------------------------------
\subsection{Remarks on the formal statements}\label{pf:sh-hodge:remarks}

\begin{enumerate}
\item Several hypotheses of the formal statements are not used in their proofs: in \cref{pf:sh-hodge:tek} the three literature hypotheses and the
  comparison data; in \cref{pf:sh-hodge:lc} the three literature hypotheses; in \cref{pf:sh-hodge:tb} the comparison data; in \cref{pf:sh-hodge:bkC}
  the local Euler characteristic formula and the tie of the local theory; in \cref{pf:sh-hodge:bka} the goodness of the coefficient extension.
\item The admitted statement (adm4) of \cref{pf:sh-hodge:def-adm} is called \emph{Archimedean signature matching} in the development, but it is the local
  Euler characteristic formula \shhodgeLEC; likewise the fourth admitted hypothesis of the lifting theory, so named, is the clause (P4).
\item The formal statements \lref{GlobalTheory.T.blochKatoLocal_all_of_tateDuality_eulerChar} and
  \lref{GlobalTheory.BlochKatoLocalOfHyp.blochKatoLocal_allPrimes_of_blochKato} are identical.
\item The formal Chebotarev statement \cref{pf:sh-hodge:cheb} is an existence statement, weaker than the density theorem; it is proved without the
  density theorem, by Chebotarev's cyclotomic crossing and Hecke's $L$-functions.
\item The statement $\sfname{Cheb}$ used in \cref{pf:sh-hodge:js} is a hypothesis of that theorem; it is equivalent to \cref{pf:sh-hodge:cheb} and
  is derived from it in the development.
\item In \cref{pf:sh-hodge:glit}, global Artin reciprocity \shhodgeGAR{} is used only through the bound for the unramified homomorphisms of $G_E$
  by the class number $h_E$.
\item The results of \cref{pf:sh-hodge:hw,pf:sh-hodge:loc} are statements about arbitrary Hodge--Tate, Fontaine and comparison data of the local theory of a
  datum (which exist by \cref{pf:sh-hodge:exhw}); the Wach structures are axiomatised, and the clauses (P1)--(P10) are theorems about every such
  axiomatised structure arising from the data.
\end{enumerate}
